\pdfoutput=1
% arXiv preprint version, generated by paper/rewrite/build_derived.py from main_cmame.tex.
% Do not edit by hand: edit main_cmame.tex and rerun the generator.
\documentclass[11pt,a4paper]{article}

\usepackage[margin=1in]{geometry}
\usepackage[T1]{fontenc}
\usepackage{lmodern}
\usepackage{amsmath,amssymb,amsfonts,amsthm}
\usepackage{array}
\usepackage{booktabs}
\usepackage{graphicx}
\usepackage{url}
\usepackage{xcolor}
\usepackage[colorlinks=true,linkcolor=blue,citecolor=blue,urlcolor=blue]{hyperref}
\providecommand{\nolinkurl}[1]{\url{#1}}
\emergencystretch=4em

\newtheorem{assumption}{Assumption}
\newtheorem{definition}{Definition}
\newtheorem{lemma}{Lemma}
\newtheorem{theorem}{Theorem}
\newtheorem{proposition}{Proposition}
\newtheorem{corollary}{Corollary}

\newcommand{\method}{Gauss6\slash FullVA}
\newcommand{\tfe}{TFE}
\newcommand{\figpath}{.}
\newcolumntype{P}[1]{>{\raggedright\arraybackslash}p{#1}}
\renewcommand{\topfraction}{0.9}
\renewcommand{\bottomfraction}{0.7}
\renewcommand{\textfraction}{0.08}
\renewcommand{\floatpagefraction}{0.85}
\setcounter{topnumber}{3}
\setcounter{totalnumber}{4}
\newcommand{\paperfig}[3]{%
  \begin{figure}[!htb]
    \centering
    \IfFileExists{#1}{%
      \includegraphics[width=\linewidth]{#1}%
    }{%
      \fbox{\parbox{0.86\linewidth}{Missing figure: \nolinkurl{#1}}}%
    }
    \caption{#2}
    \label{#3}
  \end{figure}
}

\begin{document}

\title{A sixth-order Lie-group Gauss collocation integrator for lower-pair mechanisms with position-, velocity- and acceleration-level constraints}
\author{Jingquan Wang\\[3pt]
  \small Department of Mechanical Engineering, University of Wisconsin--Madison,\\
  \small 1513 University Avenue, Madison, WI 53706, United States\\
  \small \texttt{jingquanw@wisc.edu}}
\date{}

\maketitle

\begin{abstract}
Multibody systems with lower-pair joints in absolute coordinates are index-3 differential-algebraic
equations on a Lie group. We study \method{}, a three-stage Gauss collocation step whose stage
system enforces the joint constraints at position, velocity and acceleration level together with
the Newton--Euler balance, and reconstructs the rotational endpoint through the exponential map.
The central result is an exact algebraic identity: the lifted reduced Gauss stage, obtained by
collocating the joint-coordinate equations of motion and lifting the result to absolute
coordinates, satisfies every one of the $132$ implemented stage rows, and conversely every root of
the non-dynamic rows on the regular branch is such a lift. The method is therefore reduced Gauss
collocation executed in absolute coordinates. Everything the implementation adds, the
velocity-level endpoint closure and the inexact Newton solve, is an $O(h^7)$ perturbation, which
gives a sixth-order error bound whose only standing hypothesis is smoothness on a compact
neighbourhood of the trajectory; the solver and inverse interfaces are derived for $h\le h_0$, and
$h_0$ is quantified and traced to the curvature of the friction law. The row identities, the
perturbation chain and the Butcher conditions are machine-checked in Lean~4 with Mathlib
($57$ theorems). Numerically the method is sixth order on a frictional two-body chain and on the
ASME double pendulum, reproduces the driven ASME mechanisms to roundoff, is two to seven orders of
magnitude more accurate than the one- and two-stage members of the same family at the same Newton
work, and reaches errors below $10^{-10}$ on the double pendulum where the public
absolute-coordinate codes have errors of order one. A sweep of the Stribeck velocity locates the
regime in which a sharp friction law reduces the observed order on coarse grids.
\end{abstract}

\medskip\noindent\textbf{Keywords:} Lie group integrator ; index-3 DAE ; multibody dynamics ; lower-pair mechanisms ; Gauss collocation ; machine-checked proof

\section{Introduction}
\label{sec:intro}

Multibody systems with lower-pair joints are most conveniently modelled in absolute coordinates:
each body carries its own position and rotation, and every joint contributes algebraic constraints.
The price is an index-3 differential-algebraic system on a Lie group, and the integrator has to keep
three things consistent at once: the constraints at position level, their time derivatives at
velocity and acceleration level, and the rotation-group structure of the configuration. Methods
that enforce only the position constraints let the velocity and acceleration constraints drift;
methods that stabilize them by projection or index reduction change the equations that are being
solved; and methods that integrate a reduced set of joint coordinates require the reduced equations
of motion, which are mechanism-specific to derive and lose the direct access to reaction forces
that the absolute formulation provides.

Time finite elements on Lie groups~\cite{chaturvedi2026tfe} have recently shown that high order is
attainable for frictional index-3 pendulum problems, and the absolute-coordinate formulations of
Negrut and co-workers~\cite{taves2020exponential,kissel2021dwelling,kissel2022absolute,fang2023halfimplicit}
have established a public benchmark suite of four mechanisms on which such methods can be compared.
What has been missing is a method of order higher than four for the full index-3 problem in
absolute coordinates whose order can be proved rather than observed, and whose proof survives the
details that implementations add: acceleration-level constraint rows, Lagrange multipliers as stage
unknowns, a Newton solver with a finite tolerance, and an endpoint reconstruction through the
exponential map.

This paper studies such a method, \method{}. Its one-step map solves a single square nonlinear
system in which the three-stage Gauss collocation conditions for the joint coordinates, the joint
constraints at position, velocity and acceleration level, and the Newton--Euler balance at every
stage appear as rows (Section~\ref{sec:method}); the rotational endpoint is then reconstructed from
the Gauss quadrature of the body angular velocities through the exponential map. The stage system
has $132$ rows for the two-body chain used throughout, and the same construction applies to the
four benchmark mechanisms up to the joint list.

The central result is algebraic. We show that the lifted reduced Gauss stage, that is, the stage
vector obtained by taking the three-stage Gauss collocation solution of the reduced joint-coordinate
equations of motion and lifting it to absolute coordinates, makes every one of the $132$ rows
vanish, and conversely that any root of the non-dynamic rows on the regular branch is such a lift
(Lemma~\ref{lem:exact-stage-identity}). The implemented method is therefore reduced Gauss
collocation, executed in absolute coordinates without ever forming the reduced equations, and the
classical sixth-order theory of Gauss collocation applies to it verbatim. Everything that the
implementation adds on top of the reduced collocation step, namely the endpoint closure of the
velocity-level constraints and the inexact Newton solve, is a perturbation whose size we bound
explicitly. The resulting local error bound (Theorem~\ref{thm:g6fullva-order}) is
$C_{\rm loc}h^7 = (C_G + C_E + C_N c_\eta)h^7$, and its only standing hypothesis is the smoothness
of the constrained problem on a compact neighbourhood of the trajectory: the uniform invertibility
of the stage Jacobian and the reachability of the $c_\eta h^7$ solver tolerance follow from
smoothness for $h\le h_0$ (Lemmas~\ref{lem:p2-from-p1} and~\ref{lem:newton-envelope}). We quantify
$h_0$ on the benchmark and find that it is set by the curvature of the friction law, not by the
collocation structure.

The algebraic identities and the perturbation lemmas are machine-checked in Lean~4 with Mathlib:
the row identities are proved for a transcription of the implemented residual, the perturbation
chain is proved with its explicit constants, and the Gauss tableau is shown to satisfy Butcher's
simplifying assumptions $B(6)$, $C(3)$, $D(3)$ and to fail $B(7)$. Butcher's order theorem and the
classical Gauss local-error estimate enter as cited results. The development ships with the paper
(\ref{app:lean}).

Numerically, the method is sixth order on a frictional two-body chain over its regular branch for
six step sizes (Section~\ref{sec:numerical}) and on the ASME double pendulum, the one benchmark
mechanism of the public suite whose motion is not prescribed by a drive; the three driven
mechanisms, including the two closed loops, are reproduced to roundoff (the driven pendulum from
$h=0.025$ on, the loops at every step size). The
sixth-order stage system is repaid in work/precision terms against the lower-order members of the
same Gauss family, which cost the same number of Newton iterations per step, and the double pendulum
is compared on the public horizon with the public absolute-coordinate codes against one common
reference. A sweep of the friction regularization locates the regime in which a sharp Stribeck law
reduces the observed order on coarse grids, and the constraint norms and the energy of the
frictionless variant are followed over the whole regular branch. What we do not do is reproduce the
temporal finite-element method of~\cite{chaturvedi2026tfe} on its own problems; the comparison with
that method is at the level of formal order and is stated as such.

\paragraph{Contributions}
\begin{enumerate}
  \item The exact stage identity: the FullVA stage system with acceleration-level constraints is
  reduced three-stage Gauss collocation lifted to absolute coordinates, in both directions.
  \item A sixth-order local and global error theorem for the implemented one-step map under a
  smoothness assumption alone, with the solver and inverse interfaces derived and the small-step
  threshold quantified.
  \item A Lean~4/Mathlib development of $57$ theorems covering the row identities, the perturbation
  chain with explicit constants, and the tableau order conditions.
  \item Convergence, work/precision, friction-regime and constraint-drift results on a frictional
  chain and on the four public benchmark mechanisms, all generated from a single implementation path
  and reproducible from the accompanying repository.
\end{enumerate}

Section~\ref{sec:related} reviews Lie-group and index-3 integrators for multibody systems.
Section~\ref{sec:setting} fixes the setting and notation. Section~\ref{sec:method} defines the
one-step map. Section~\ref{sec:order} contains the identity, the perturbation lemmas and the
theorem. Section~\ref{sec:numerical} reports the experiments, Section~\ref{sec:limitations} the
limitations, and Section~\ref{sec:conclusions} concludes.

\section{Related work}
\label{sec:related}

\paragraph{Lie-group integrators}
Integrators that advance rotations in a local tangent chart and return to the group through the
exponential map go back to Crouch and Grossman~\cite{crouch1993numerical} and to the
Runge--Kutta--Munthe-Kaas construction~\cite{munthe1998lie,munthe1999high}; the surveys of
Iserles et al.~\cite{iserles2000lie} and of Celledoni, Marthinsen and Owren~\cite{celledoni2014lie}
and the monograph of Hairer, Lubich and Wanner~\cite{hairer2006geometric} give the general theory.
In multibody dynamics the same viewpoint underlies the Lie-group generalized-$\alpha$ schemes of
Br\"uls, Cardona and Arnold~\cite{bruls2010lie,bruls2012lie}, whose convergence for constrained
systems was analysed in~\cite{arnold2007convergence,arnold2015error}, BDF schemes on Lie
groups~\cite{wieloch2021bdf}, the geometrically exact beam and rigid-body formulations on
$SE(3)$~\cite{sonneville2014geometrically}, and singularity-free quaternion
integration~\cite{terze2016singularity,terze2020aircraft}. Bottasso and Borri~\cite{bottasso1998integrating}
discuss the interaction between the rotation parametrization and the integrator. These methods
establish how to evolve rotations geometrically; they are of order at most two, or, in the
Runge--Kutta--Munthe-Kaas case, are formulated for ordinary differential equations on the group
rather than for the constrained index-3 problem.

\paragraph{Constrained mechanical systems and index-3 DAEs}
Absolute-coordinate multibody systems are index-3 differential-algebraic equations
\cite{hairer1993solving,ascher1998computer,petzold1986numerical}. Their numerical order depends
not only on the integrator but on how the constraints and their differentiated forms are enforced:
index reduction with stabilization~\cite{gear1985automatic}, projection, and the review of
constraint-enforcement approaches by Bauchau and Laulusa~\cite{bauchau2008review} describe the
options and their failure modes. Runge--Kutta methods for index-2 and overdetermined index-3
formulations were analysed by Jay~\cite{jay1996symplectic,jay2006specialized,jay2007specialized};
the Lobatto IIIA--IIIB pairs of that work are the closest relatives of the present stage system in
that they enforce the position and velocity constraints at every stage. Lunk and
Simeon~\cite{lunk2006solving} and Arnold and Br\"uls~\cite{arnold2007convergence} treat the
Newmark and generalized-$\alpha$ families for constrained systems. The absolute-coordinate
formulations of Negrut and co-workers
\cite{taves2020exponential,kissel2021dwelling,kissel2022absolute,fang2023halfimplicit,kissel2024velocitypartitioning}
compare three rotation parametrizations and a half-implicit scheme on a public suite of four
mechanisms; that suite and its public code are used in Section~\ref{sec:numerical}. Reduced
(joint-coordinate) formulations, for which the constrained problem becomes an ordinary
differential equation, are classical~\cite{nikravesh1988computer,shabana2020dynamics,jain2011robot};
the present paper shows that a particular absolute-coordinate collocation step coincides with
Gauss collocation of such a reduced formulation without forming it.

\paragraph{Collocation and variational integrators}
Gauss collocation is the standard high-order implicit Runge--Kutta family: the $s$-stage method has
order $2s$, is symplectic and $B$-stable, and its stages are the values of a collocation polynomial
\cite{hairer1993nonstiff,butcher2016numerical,hairer2006geometric}. Variational and symplectic
integrators~\cite{marsden2001discrete,leyendecker2008variational,leok2012prolongation,bourabee2009hamilton}
derive the discrete equations from a discrete action and treat holonomic constraints at the
discrete level; Betsch and Steinmann~\cite{betsch2002conservation} obtain energy--momentum
conserving time finite elements for constrained systems. The present method is not variational: it
enforces the constraint hierarchy pointwise at the stages and leaves the conservation properties
to those of Gauss collocation on the reduced problem.

\paragraph{Friction and nonsmooth dynamics}
Regularized friction laws such as the Brown--McPhee model~\cite{browmcphee2016friction} replace the
set-valued Coulomb law by a smooth velocity-dependent force; Pennestr\`{\i} et
al.~\cite{pennestri2016review} review the alternatives. Nonsmooth time-stepping and complementarity
formulations~\cite{stewart1996implicit,anitescu1997formulating,tasora2011matrixfree} handle
unilateral contact and Coulomb friction without regularization but are of low order by design. The
problems in this paper are bilateral lower-pair mechanisms with regularized friction; the effect of
the regularization parameter on the observed order is quantified in Section~\ref{sec:numerical}.

\paragraph{Time finite elements}
Time finite elements for structural and rigid-body dynamics~\cite{borri1988rigid,hughes1988spacetime,hulbert1992time,kim2017higher,jog2016robust,bauchau2024time}
provide a different route to high order. The Lie-group time finite-element method of Chaturvedi,
Sandu and Sandu~\cite{chaturvedi2026tfe} extends this route to frictional index-3 multibody DAEs;
its Gauss--Lobatto family of degree $m$ has formal order $2m-1$, so that its $m=3$ member is a
fifth-order formula, although the source paper reports order reduction on the DAE problem. We do
not reimplement that method; the comparison with it in this paper is at the level of formal order.

\section{Problem setting}
\label{sec:setting}

\paragraph{Equations of motion}
A rigid body $b$ has centre-of-mass position $r_b\in\mathbb R^3$, rotation $Q_b\in SO(3)$,
translational velocity $v_b=\dot r_b$, body-frame angular velocity $\omega_b$ with
$\dot Q_b=Q_b\,\widehat{\omega_b}$, mass $m_b$ and body-frame inertia $J_b$. With
$q=(r_b,Q_b)_b$, $v=(v_b,\omega_b)_b$ and $a=\dot v$, a mechanism with lower-pair joints is
governed by
\begin{align}
  M\dot v + G(v) + \Phi_q(q)^T\lambda &= F(q,v,t), \label{eq:dae-dyn}\\
  \dot q &= \mathcal K(q)v, \label{eq:dae-kin}\\
  \Phi(q) &= 0, \label{eq:dae-constraint}
\end{align}
where $M=\operatorname{diag}(m_bI,J_b)$, $G$ collects the gyroscopic terms
$\omega_b\times J_b\omega_b$, $\Phi$ stacks the joint constraints, $\lambda$ are the Lagrange
multipliers, and $\mathcal K$ is the kinematic map, which is the identity on the translational part
and $\omega_b\mapsto Q_b\widehat{\omega_b}$ on the rotational part. The system is an index-3
differential-algebraic equation on $(\mathbb R^3\times SO(3))^{N}$. Differentiating the constraint
once and twice gives the velocity- and acceleration-level constraints
\begin{equation}
  \Phi_q(q)v=0,\qquad \Phi_q(q)a+\dot\Phi_q(q,v)v=0,
  \label{eq:fullva}
\end{equation}
which the exact solution satisfies identically; the method of Section~\ref{sec:method} enforces
Eq.~\eqref{eq:dae-constraint} and both equations of Eq.~\eqref{eq:fullva} at every stage, which
is what the acronym FullVA (full velocity- and acceleration-level enforcement) refers to.

\paragraph{Rotations}
Rotations are represented by unit quaternions in the implementation and advanced by the retraction
\begin{equation}
  Q_{n+1}=Q_n\,\mathrm{Exp}(\eta_{n+1}),\qquad Q_i=Q_n\,\mathrm{Exp}(\eta_i),
  \label{eq:retraction}
\end{equation}
with $\eta\in\mathbb R^3\cong\mathfrak{so}(3)$ the local rotation increment and $\mathrm{Exp}$
the exponential map. The right-trivialized differential of the exponential map is written
$J_r(\eta)$, so that $\dot\eta=J_r^{-1}(\eta)\,\omega$ along a motion with body angular velocity
$\omega$.

\paragraph{Lower pairs and joint coordinates}
A lower pair $j$ connects a parent body $p(j)$ (possibly the ground) and a child body $c(j)$. Let
$P_j(q)$ and $C_j(q)$ be the attachment points on the two bodies, $d_j(q)=C_j(q)-P_j(q)$ the
relative displacement, $a_j(q)$ the joint axis carried by the parent, $\tilde a_j(q)$ the axis
carried by the child, and $\omega^{\rm rel}_j=Q_{c(j)}\omega_{c(j)}-Q_{p(j)}\omega_{p(j)}$ the
relative angular velocity in the world frame. Each pair carries a fixed world direction $n_j$, the
joint axis at $t=0$, and a fixed orthonormal pair $e_{j,1},e_{j,2}\perp n_j$. A pair with two
degrees of freedom (sliding along and rotating about its axis) has the joint coordinates
\begin{equation}
  s_j(q)=n_j\cdot d_j(q),\qquad
  \theta_j(q)=\operatorname{atan2}\bigl(\tilde t_j\cdot(a_j\times t_j),\ \tilde t_j\cdot t_j\bigr),
  \label{eq:joint-coordinates}
\end{equation}
where $t_j,\tilde t_j$ are twist reference vectors fixed in the parent and in the child. Their
rates are $\dot s_j=n_j\cdot\dot d_j$ and $\dot\theta_j=a_j\cdot\omega^{\rm rel}_j$, and their
second rates are $\ddot s_j=n_j\cdot\ddot d_j$ and
$\ddot\theta_j=a_j\cdot\dot\omega^{\rm rel}_j+\dot a_j\cdot\omega^{\rm rel}_j$. Revolute,
prismatic and spherical joints are obtained by fixing one or both coordinates or by dropping the
alignment rows; the constraint rows of the four benchmark mechanisms are listed in
\ref{app:lower-pair}.

\paragraph{The two-body chain}
The main benchmark is a two-body open chain (Table~\ref{tab:chain}). Body $0$ is connected to the
ground by a pair with sliding and rotation along a fixed ground axis; body $1$ is connected to body
$0$ by a second pair whose axis $a_1(q)=Q_0\hat a^{\rm next}_0$ is fixed in body $0$. The two axes
fixed in body $0$ are not parallel and body $0$ spins, so $a_1(q)$ moves on a cone about the ground
axis while $n_1$ and $e_{1,m}$ stay frozen at their $t=0$ values. The second pair is therefore a
four-constraint lower pair with a fixed sliding direction and a moving rotation axis, not a
cylindrical pair in the strict sense; its friction sliding speed is $\dot s_1\,(n_1\cdot a_1)$, and
the modelling is regular as long as $n_1\cdot a_1(q)\neq0$. On the initial state of
Table~\ref{tab:chain} this factor decreases from $1$ to $0.34$ on $[0,0.5]$ and vanishes at
$t\approx0.605$, where the frozen sliding direction ceases to define the pair; all chain
experiments are run on the regular branch $[0,0.5]$. Each joint carries gravity, constant external
loads, viscous damping and a regularized Brown--McPhee friction torque and force
\cite{browmcphee2016friction} with static and dynamic coefficients $\mu_s,\mu_d$ and Stribeck
velocity $v_s$; the friction force enters the Newton--Euler rows through the multipliers (normal
force) and the sliding speed, and is $C^\infty$ in the state for $v_s>0$. The smooth case uses
$v_s=0.5$; Section~\ref{sec:numerical} varies $v_s$ down to $0.02$.

\begin{table}[!htb]
\centering
\caption{Two-body chain benchmark. Directions are given before normalization; the initial rotations
are the unique rotations mapping the body-fixed axis and twist vectors of each pair onto their
parent-side counterparts.}
\label{tab:chain}
\small
\begin{tabular}{P{0.42\linewidth}P{0.52\linewidth}}
\toprule
Quantity & Value \\
\midrule
Masses $m_0,m_1$ [kg] & $2.7$, $1.9$ \\
Inertias $J_0,J_1$ [kg\,m$^2$] & $\operatorname{diag}(0.11,0.24,0.31)$, $\operatorname{diag}(0.08,0.19,0.27)$ \\
Ground axis, ground twist & $(0.31,0.78,0.54)$, first vector of an orthonormal complement \\
Parent-side axes $\hat a^{\rm prev}_0,\hat a^{\rm prev}_1$ & $(0.18,0.91,-0.37)$, $(-0.24,0.87,0.43)$ \\
Child-side axes $\hat a^{\rm next}_0,\hat a^{\rm next}_1$ & $(0.42,0.69,0.59)$, $(0,1,0)$ \\
Attachment points $s^{\rm prev}_0,s^{\rm prev}_1$ [m] & $(0.12,-0.04,0.18)$, $(-0.08,0.05,0.16)$ \\
Attachment points $s^{\rm next}_0,s^{\rm next}_1$ [m] & $(0.16,0.02,-0.20)$, $(0,0,0)$ \\
Gravity [m/s$^2$] & $(0,0,-9.81)$ \\
External forces (world) [N] & $(0.8,-0.35,0.25)$, $(-0.25,0.45,-0.10)$ \\
External torques (body) [N\,m] & $(0.03,-0.02,0.015)$, $(-0.018,0.026,-0.014)$ \\
Friction $\mu_s$, $\mu_d$, viscous coefficient, radius & $0.30$, $0.20$, $0.030$, $1.0$ \\
Stribeck velocity $v_s$ & $0.5$ (smooth case); $0.2,0.1,0.05,0.02$ in the sweep \\
Initial $s_j(0)$, $\dot s_j(0)$, $\dot\theta_j(0)$ & $(0.35,0.42)$, $(0.42,-0.26)$, $(0.18,-0.14)$ \\
\bottomrule
\end{tabular}
\end{table}

\paperfig{\figpath/Figure_1_chain_schematic.png}
{The two-body chain: body $0$ slides and spins along the fixed ground axis, body $1$ slides and spins
along the axis $a_1(q)$ carried by body $0$; the second pair keeps its sliding direction $n_1$ and
its basis $e_{1,m}$ frozen at their $t=0$ values.}
{fig:chain}

\paragraph{Notation}
Throughout, $h$ is the step size, $t_n=nh$, $x_n=(q_n,v_n)$ the endpoint state (accelerations and
multipliers are stage quantities), $c_i,b_i,A_{ik}$ ($i,k=1,2,3$) the nodes, weights and matrix of
the three-stage Gauss tableau, $Z$ the stacked stage vector, and $\|\cdot\|$ a fixed norm on the
finite-dimensional space at hand. Constants $C$ with subscripts are independent of $h$ and of the
number of steps.

\section{The \method{} step}
\label{sec:method}

\subsection{Stage unknowns and stage system}

For each stage $i=1,2,3$ and each body $b$ the unknowns are the translational position $r_{b,i}$,
the rotation increment $\eta_{b,i}$ with $Q_{b,i}=Q_{b,n}\mathrm{Exp}(\eta_{b,i})$, the velocities
$v_{b,i},\omega_{b,i}$ and the accelerations $a_{b,i},\alpha_{b,i}$; each lower pair $j$
contributes its multipliers $\lambda_{j,i}$, four per pair for the chain. With two bodies and two
pairs the stage vector
\begin{equation}
  Z=(Z_1,Z_2,Z_3),\qquad Z_i=(r_i,\eta_i,v_i,\omega_i,a_i,\alpha_i,\lambda_i),
  \label{eq:stage-variable-block}
\end{equation}
has $3\,(2\cdot18+2\cdot4)=132$ components.

Write $\psi_j(q)=\tilde a_j(q)-a_j(q)$ for the axis-alignment residual and, for a joint-coordinate
function $g$ with rate $\dot g$,
\begin{equation}
  \Delta_i[g]=g(Z_i)-g(x_n)-h\sum_{k=1}^3A_{ik}\,\dot g(Z_k)
  \label{eq:collocation-defect-operator}
\end{equation}
for its Gauss collocation defect at stage $i$. With $t_i=t_n+c_ih$, the applied wrench
$(f_{b,i},\tau_{b,i})$ evaluated at the stage, and the total force and torque
$F^r_{b,i}=f_{b,i}+\Phi_{r,b}(q_i)^T\lambda_i$, $T^\theta_{b,i}=\tau_{b,i}+\Phi_{\theta,b}(q_i)^T\lambda_i$
including the multiplier wrench, the stage system reads, for every pair $j$, every $m=1,2$ and
every body $b$,
\begin{equation}
\label{eq:g6fullva-expanded-stage}
\begin{aligned}
0 &= e_{j,m}\cdot d_j(q_i), &
0 &= e_{j,m}\cdot\psi_j(q_i),\\
0 &= e_{j,m}\cdot\dot d_j(q_i,v_i), &
0 &= e_{j,m}\cdot\dot\psi_j(q_i,v_i),\\
0 &= e_{j,m}\cdot\ddot d_j(q_i,v_i,a_i), &
0 &= e_{j,m}\cdot\ddot\psi_j(q_i,v_i,a_i),\\
0 &= \Delta_i[s_j], & 0 &= \Delta_i[\dot s_j],\\
0 &= \Delta_i[\theta_j], & 0 &= \Delta_i[\dot\theta_j],\\
0 &= m_b a_{b,i}-F^r_{b,i}, &
0 &= J_b\alpha_{b,i}+\omega_{b,i}\times J_b\omega_{b,i}-T^\theta_{b,i} .
\end{aligned}
\end{equation}
The first three lines are the lower-pair constraints at position, velocity and acceleration level
($12$ rows per pair and stage); the fourth and fifth lines are Gauss collocation of the four joint
coordinates $(s_j,\dot s_j,\theta_j,\dot\theta_j)$ ($4$ rows per pair and stage); the last line is
the Newton--Euler balance ($6$ rows per body and stage). Table~\ref{tab:rows} gives the inventory:
$44$ rows per stage for the $44$ unknowns, $72+24+36=132$ over the three stages. The absolute body
coordinates $(r_i,\eta_i,v_i,\omega_i)$ carry no collocation row of their own; they are determined
by the constraint rows once the joint coordinates are collocated. The implementation projects the
two translational collocation rows onto the current joint axis $a_j(q_i)$ rather than onto $n_j$;
on the constraint manifold $d_j,\dot d_j,\ddot d_j\parallel n_j$, so the two forms differ by the
factor $n_j\cdot a_j(q_i)$, nonzero on the regular branch, and define the same stage solution.
Driven mechanisms add prescribed right-hand sides to the three constraint levels of the driven
row. The multiplier wrench collects the joint normal forces, the axis-alignment torques and the
friction force along the joint axis; the friction law enters only there.

\begin{table}[!htb]
\centering
\caption{Row inventory of the stage system for the two-body chain (per stage and in total).}
\label{tab:rows}
\small
\begin{tabular}{P{0.46\linewidth}P{0.16\linewidth}P{0.10\linewidth}P{0.10\linewidth}}
\toprule
Row family & Per pair or body & Per stage & Total \\
\midrule
Position-level constraints $\Phi=0$ & $4$ per pair & $8$ & $24$ \\
Velocity-level constraints $\Phi_qv=0$ & $4$ per pair & $8$ & $24$ \\
Acceleration-level constraints $\Phi_qa+\dot\Phi_qv=0$ & $4$ per pair & $8$ & $24$ \\
Joint-coordinate collocation $\Delta_i[s_j],\Delta_i[\dot s_j],\Delta_i[\theta_j],\Delta_i[\dot\theta_j]$ & $4$ per pair & $8$ & $24$ \\
Newton--Euler balance & $6$ per body & $12$ & $36$ \\
\midrule
Unknowns $(r,\eta,v,\omega,a,\alpha)$ per body, $\lambda$ per pair & $18$, $4$ & $44$ & $132$ \\
\bottomrule
\end{tabular}
\end{table}

\paperfig{\figpath/Figure_2_method_overview.png}
{One step of \method{}: the stage system couples the three constraint levels, the joint-coordinate
collocation rows and the Newton--Euler balance; the endpoint is reconstructed with the Gauss weights
and the exponential map and closed at velocity level.}
{fig:overview}

\subsection{Endpoint reconstruction and closure}

After the stage solve the endpoint is reconstructed with the Gauss weights,
\begin{align}
  r_{n+1} &= r_n+h\sum_{i=1}^3 b_i v_i, &
  v_{n+1} &= v_n+h\sum_{i=1}^3 b_i a_i, \label{eq:end-trans}\\
  Q_{n+1} &= Q_n\mathrm{Exp}\!\left(h\sum_{i=1}^3 b_i J_r^{-1}(\eta_i)\omega_i\right), &
  \omega_{n+1} &= \omega_n+h\sum_{i=1}^3 b_i \alpha_i .
  \label{eq:end-rot}
\end{align}
The reconstructed endpoint satisfies the position-level constraints only up to the quadrature
error ($O(h^7)$ per step, Lemma~\ref{lem:gauss-endpoint-defect}) and the velocity-level constraints
likewise. The endpoint closure map $\mathcal C_h$ then projects the endpoint velocities
$(v_{n+1},\omega_{n+1})$ onto the velocity-level constraint manifold of Eq.~\eqref{eq:fullva} at
$q_{n+1}$ by the minimal-norm correction, that is, by solving the KKT system with matrix
$\bigl[\begin{smallmatrix}I&\Phi_q^T\\ \Phi_q&0\end{smallmatrix}\bigr]$; positions and rotations are
not projected. The correction is an $O(h^7)$ perturbation of the step (Lemma~\ref{lem:endpoint-closure}).
The closed endpoint $x_{n+1}=\mathcal C_h(\mathcal E_h(Z))$, with $\mathcal E_h$ the map of
Eqs.~\eqref{eq:end-trans}--\eqref{eq:end-rot}, defines the one-step map
$\Psi_h(x_n)=x_{n+1}$.

\subsection{Newton solve and predictor}

The stage system $F_h(Z;x_n)=0$ is solved by Newton's method,
\begin{equation}
  D_ZF_h(Z^{(k)};x_n)\,\Delta Z=-F_h(Z^{(k)};x_n),\qquad Z^{(k+1)}=Z^{(k)}+\Delta Z,
  \label{eq:newton}
\end{equation}
with the $132\times132$ Jacobian assembled by forward-mode automatic differentiation of the single
residual function and factorized densely. The iteration stops when the residual norm or the update norm falls below the tolerance
$\tau_N=10^{-11}$ and is declared failed after $k_{\max}=80$ iterations. The predictor
extrapolates positions and velocities from the endpoint state with the constant endpoint
acceleration, sets the rotation increments, angular velocities and angular accelerations of the
stages to zero, and estimates the multipliers from the endpoint force balance; the consequences of
this choice for the convergence analysis are discussed after Lemma~\ref{lem:newton-envelope}. On
the chain the solve takes between five and six Newton iterations per step at every step size
(Section~\ref{sec:numerical}).

\begin{table}[!htb]
\centering
\caption{One step of \method{}.}
\label{tab:algorithm}
\small
\begin{tabular}{P{0.06\linewidth}P{0.84\linewidth}}
\toprule
Step & Operation \\
\midrule
1 & From $x_n=(q_n,v_n)$ and $h$, form the predictor $Z^{(0)}$. \\
2 & Assemble the stage residual $F_h(Z;x_n)$ of Eq.~\eqref{eq:g6fullva-expanded-stage}:
constraint rows at the three levels, joint-coordinate collocation rows, Newton--Euler rows. \\
3 & Iterate Eq.~\eqref{eq:newton} with the automatic-differentiation Jacobian until
$\|F_h\|\le\tau_N$ or $\|\Delta Z\|\le\tau_N$. \\
4 & Reconstruct the endpoint by Eqs.~\eqref{eq:end-trans}--\eqref{eq:end-rot}. \\
5 & Apply the velocity-level closure $\mathcal C_h$; record the Newton count and the endpoint
constraint norms. \\
\bottomrule
\end{tabular}
\end{table}

\paragraph{Cost}
One Newton iteration costs one residual and Jacobian evaluation and one dense $132\times132$
factorization; a step costs five to six such iterations. The Jacobian is block-sparse (each stage
couples to the others only through the collocation rows), which the present reference
implementation does not exploit. The two- and one-stage members of the same family, obtained by
replacing the Gauss tableau, have $88$ and $44$ unknowns and are used for the work/precision
comparison of Section~\ref{sec:numerical}.

\section{Error analysis}
\label{sec:order}

The analysis has three parts. Lemma~\ref{lem:exact-stage-identity} shows that the stage system of
Eq.~\eqref{eq:g6fullva-expanded-stage} is exactly three-stage Gauss collocation of the reduced
joint-coordinate equations of motion, written in absolute coordinates; the stage residual vanishes
identically at the lifted reduced Gauss stage and no stage-residual perturbation term appears.
The two perturbations the implementation does introduce, the endpoint closure and the inexact
Newton solve, are bounded by $O(h^7)$ (Lemmas~\ref{lem:endpoint-closure}
and~\ref{lem:inexact-newton}), and the local defect is transferred to the grid by a discrete
Gronwall argument (Lemma~\ref{lem:local-global-transfer}). The hypotheses are collected in
Assumption~\ref{ass:regularity}; Lemmas~\ref{lem:p2-from-p1} and~\ref{lem:newton-envelope} show that
the second and third of them follow from the first for small steps. The algebraic identities and the
perturbation lemmas are machine-checked in Lean~4/Mathlib (\ref{app:lean}); Butcher's
order theorem for the Gauss tableau and the classical Gauss local-error estimate
\cite{hairer1993solving} are used as cited results.

\subsection{Reduced chart, lift and hypotheses}

Let $y=(s_0,\theta_0,s_1,\theta_1,\dot s_0,\dot\theta_0,\dot s_1,\dot\theta_1)$ collect the joint
coordinates of Eq.~\eqref{eq:joint-coordinates} and their rates. On the smooth constrained branch
the equations of motion reduce to
\begin{equation}
  \dot y=f(y,t),
  \label{eq:reduced-ode}
\end{equation}
where $f$ returns the rates and the joint accelerations obtained from the constrained Newton--Euler
equations. Let $\varphi_h$ denote the flow of Eq.~\eqref{eq:reduced-ode} over one step,
$\mathcal L(y,t)=(q,v,a,\lambda)$ the smooth lift to absolute coordinates given by forward
kinematics together with the velocities, accelerations and multipliers of the constrained dynamics,
and $Y_1,Y_2,Y_3$ the Gauss collocation stages of Eq.~\eqref{eq:reduced-ode} from $y_n$,
\begin{equation}
  Y_i=y_n+h\sum_{k=1}^3A_{ik}f(Y_k,t_k),\qquad i=1,2,3 .
  \label{eq:reduced-gauss-stages}
\end{equation}
The lifted reduced Gauss stage is $Z_G=(\mathcal L(Y_1),\mathcal L(Y_2),\mathcal L(Y_3))$. We write
$F_h(\cdot;x_n)$ for the implemented $132$-row residual, $\mathcal E_h$ for the endpoint
reconstruction of Eqs.~\eqref{eq:end-trans}--\eqref{eq:end-rot}, $\mathcal C_h$ for the endpoint
closure, $\mathcal P_h=\mathcal C_h\circ\mathcal E_h$, and $\Psi_h^G=\mathcal E_h(Z_G)$ for the
reconstructed Gauss endpoint. Norms are fixed finite-dimensional norms on the reduced chart, the
stage space and the endpoint space; chart-to-configuration norm equivalences on the compact tube are
absorbed into the constants.

\begin{assumption}[Regularity and perturbation scale]
\label{ass:regularity}
Fix a final time $T$ and a compact reduced-state tube $K$ containing the exact trajectory of
Eq.~\eqref{eq:reduced-ode} at distance at least $d_K>0$ from $\partial K$. The following hold
uniformly for $y\in K$, $0\le t\le T$ and $0<h\le h_0$, with constants independent of $h$ and of
the number of steps.
\begin{enumerate}
  \item[(H1)] (Smooth branch.) The constraint Jacobian $\Phi_q$ has constant full row rank, the
  constrained mass matrix is nonsingular on the tangent space, $f$ is $C^7$, the lift $\mathcal L$
  is smooth on the tube, and the joint coordinates complete the constraints to a chart, i.e.\ the
  map $q\mapsto(\Phi(q),s_0(q),\theta_0(q),s_1(q),\theta_1(q))$ has an invertible Jacobian.
  \item[(H2)] (Local inverses and stability.) The stage Jacobian $J_h=D_ZF_h(Z_G;x_n)$ is
  invertible with $\|J_h^{-1}\|\le M$ and $\|D_Z^2F_h\|\le L$ on a ball $B_\rho(Z_G)$ with
  $ML\rho\le\tfrac12$. The closure subsystem $E$ of Lemma~\ref{lem:endpoint-closure} is $C^2$ on a
  ball $B_{r_E}(z_\ast)$ around a same-branch closed anchor $z_\ast$ with $E(z_\ast)=0$;
  $DE(z_\ast)$ has a right inverse $B_\ast$ with $\|B_\ast\|\le M_{E}$,
  $\|DE(z)-DE(z_\ast)\|\le L_E\|z-z_\ast\|$ on that ball, $M_{E}L_Er_E\le\tfrac12$, the unclosed
  endpoint $z_u=\Psi_h^G$ lies in $B_{r_E/2}(z_\ast)$, and its raw closure defect satisfies
  $\|E(z_u)\|\le C_{E,\mathrm{raw}}h^7$. The one-step map satisfies
  \begin{equation}
  \label{eq:assumption-stability-scale}
    \|\Psi_h(y)-\Psi_h(\bar y)\|\le (1+C_sh)\|y-\bar y\|,\qquad y,\bar y\in K .
  \end{equation}
  \item[(H3)] (Solver envelope.) The computed Newton iterate $\widetilde Z$ lies in the ball
  $B_\rho(Z_G)$ of (H2) and its residual satisfies $\|F_h(\widetilde Z;x_n)\|\le\eta_h\le c_\eta h^7$.
  The output map satisfies $\|D_Z\mathcal P_h\|\le M_N$ on that ball.
\end{enumerate}
\end{assumption}

(H1) is the compact-tube regularity hypothesis of classical collocation theory. (H2) and (H3) are
the interfaces through which the implementation enters; Lemmas~\ref{lem:p2-from-p1}
and~\ref{lem:newton-envelope} derive them from (H1) for $h\le h_0$, and they are kept as separate
items only because the theorem is stated for any mechanism whose stage system has the form of
Eq.~\eqref{eq:g6fullva-expanded-stage}. (In the accompanying repository the three hypotheses are
called P1, P2 and P6.) The theorem also uses the classical Gauss local-error estimate: for
$f\in C^7$ on the compact tube there is $C_{G,0}$ with
\begin{equation}
  \|u(t_n+h)-\varphi_h(y_n)\|\le C_{G,0}h^7,
  \label{eq:gauss-local-truncation}
\end{equation}
where $u$ is the degree-three collocation polynomial through Eq.~\eqref{eq:reduced-gauss-stages}
\cite{hairer1993solving}. The tableau used in the implementation satisfies Butcher's simplifying
assumptions $B(6)$, $C(3)$, $D(3)$ and fails $B(7)$, so its order is exactly six; these algebraic
conditions are machine-checked for the coefficients hard-coded in the implementation.

\subsection{Exact stage identity}

\begin{lemma}[Exact stage identity]
\label{lem:exact-stage-identity}
Under (H1), the implemented $132$-row residual satisfies
\begin{equation}
  F_h(Z_G;x_n)=0 .
  \label{eq:exact-stage-identity}
\end{equation}
Conversely, on the branch where $n_j\cdot a_j(q_i)\neq0$, every stage vector that satisfies the
$96$ non-dynamic rows of Eq.~\eqref{eq:g6fullva-expanded-stage} lies on the constraint manifold
with consistent velocities and accelerations and has joint coordinates satisfying
Eq.~\eqref{eq:reduced-gauss-stages}. Hence the stage solve is three-stage Gauss collocation of
Eq.~\eqref{eq:reduced-ode} in the joint-coordinate chart, written in absolute coordinates.
\end{lemma}

\begin{proof}
The $72$ constraint rows vanish on $Z_G$ because each $\mathcal L(Y_k)$ lies on the constraint
manifold with consistent velocities and accelerations. The $24$ joint-coordinate rows are, by
Eq.~\eqref{eq:joint-coordinates}, the reduced collocation equations~\eqref{eq:reduced-gauss-stages}
read component-wise, because $s_j,\theta_j,\dot s_j,\dot\theta_j$ are coordinates of $y$ and
$\ddot s_j,\ddot\theta_j$ are components of $f$; the implemented projection onto $a_j(q_i)$
multiplies each translational row by $n_j\cdot a_j(q_i)$. The $36$ Newton--Euler rows vanish because
$\mathcal L$ supplies the stage accelerations and multipliers satisfying force and moment balance,
and the implemented rows are exactly these balance defects with the implemented wrench sign
conventions. For the converse, the constraint rows force $d_j,\dot d_j,\ddot d_j\parallel n_j$, so
each translational collocation row factors as
$\bigl(s_j(Z_i)-s_j(x_n)-h\sum_kA_{ik}\dot s_j(Z_k)\bigr)\,(n_j\cdot a_j(q_i))$, and the twist rows
are already scalar. Both directions are verified row by row for a transcription of the implemented
residual in the Lean development.
\end{proof}

The lemma replaces any stage-residual perturbation estimate: the residual at the lifted Gauss
stage is not small, it is zero. What remains to be controlled is which root Newton selects, how the
endpoint is reconstructed and closed, and how inexact the Newton solve is. The lemma is stated for
the open chain, where the joint coordinates form a global chart; for closed loops the same argument
applies on any local chart of independent coordinates. The rows of
Eq.~\eqref{eq:g6fullva-expanded-stage} are transcribed from the single implemented residual
function, whose automatic derivative is the Newton Jacobian, so there is one residual and Jacobian
path; the Lean development transcribes the same rows, and the binding between transcription and
source is a static check, not a theorem. Numerically, the converged stage vectors of the chain runs
read in the joint-coordinate chart have reduced collocation defects and out-of-axis components of
$d_j,\dot d_j,\ddot d_j$ at roundoff (Section~\ref{sec:numerical}).

\begin{lemma}[Branch selection]
\label{lem:stage-residual-defect}
Under (H2), $Z_G$ is the unique zero of $F_h(\cdot;x_n)$ in $B_\rho(Z_G)$. In particular the exact
stage root on that branch equals $Z_G$, and $\mathcal E_h(Z_G)=\Psi_h^G$.
\end{lemma}

\begin{proof}
For $Z\in B_\rho(Z_G)$ the mean value inequality gives
$\|F_h(Z)-F_h(Z_G)-J_h(Z-Z_G)\|\le\tfrac{L\rho}{2}\|Z-Z_G\|$. If $F_h(Z)=0$, then
Eq.~\eqref{eq:exact-stage-identity} gives
$\|Z-Z_G\|\le M\|F_h(Z)-F_h(Z_G)-J_h(Z-Z_G)\|\le\tfrac{ML\rho}{2}\|Z-Z_G\|\le\tfrac14\|Z-Z_G\|$,
so $Z=Z_G$.
\end{proof}

\begin{lemma}[(H2) from (H1) for small steps]
\label{lem:p2-from-p1}
Assume (H1). Then there is $h_0>0$ such that for $0<h\le h_0$: (i) the stage Jacobian
$J_h=D_ZF_h(Z_G)$ is invertible with
\begin{equation}
  \|J_h^{-1}\|\le 2\|J_0^{-1}\|,\qquad J_0=\lim_{h\to0}J_h,\qquad \|J_h-J_0\|\le C_Jh ,
  \label{eq:p2-perturbation-bound}
\end{equation}
(ii) the closure Jacobian is invertible, so $B_\ast$ exists with $\|B_\ast\|$ bounded uniformly,
(iii) the raw closure defect satisfies $\|E(z_u)\|\le C_{E,\mathrm{raw}}h^7$, and (iv) the one-step
map satisfies Eq.~\eqref{eq:assumption-stability-scale}. Hence all of (H2) follows from (H1) for
small steps.
\end{lemma}

\begin{proof}
At $h=0$ the collocation rows reduce to $s_j(Z_i)-s_j(x_n)$, $\theta_j(Z_i)-\theta_j(x_n)$ and
their rate analogues, so $J_0$ is block diagonal over the stages and block triangular within a
stage: the position block is the Jacobian of $q\mapsto(\Phi,s,\theta)$, the velocity block is the
same matrix acting on $(v,\omega)$ because on the constraint manifold the rate rows are the time
derivatives of the coordinate rows, and the acceleration block is the constrained Newton--Euler
system $\bigl[\begin{smallmatrix}M&\Phi_q^T\\ \Phi_q&0\end{smallmatrix}\bigr]$ acting on
$(a,\alpha,\lambda)$. Each block is invertible under (H1), so $J_0$ is invertible with an
$h$-independent inverse bound. The only $h$-dependence of $F_h$ is the factor $h\sum_kA_{ik}$ in
the collocation rows, so $\|J_h-J_0\|\le C_Jh$ on the compact tube, and for
$2\|J_0^{-1}\|C_Jh\le1$ Eq.~\eqref{eq:p2-perturbation-bound} follows from the Neumann series. The
endpoint closure is a square KKT system in the endpoint velocities whose matrix is
$\bigl[\begin{smallmatrix}I&\Phi_q^T\\ \Phi_q&0\end{smallmatrix}\bigr]$, invertible by full row rank
of $\Phi_q$; its raw defect is the velocity-level constraint value at the Gauss-weight endpoint,
which is the quadrature error of Lemma~\ref{lem:gauss-endpoint-defect} and hence $O(h^7)$; the
implicit function theorem gives a same-branch closed anchor $z_\ast$ within $O(h^7)$ of $z_u$, so
$z_u\in B_{r_E/2}(z_\ast)$ for small $h$. Finally, by Lemma~\ref{lem:exact-stage-identity} the
one-step map read in the reduced chart is the reduced Gauss step, whose derivative is $I+O(h)$ by
the implicit function theorem for the stage equations, composed with the smooth endpoint
reconstruction and an $O(h^7)$ closure correction; this gives the Lipschitz scale $1+C_sh$.
\end{proof}

The threshold $h_0$ of the Neumann argument is not sharp. On the chain
(Table~\ref{tab:p2-constants-check} in Section~\ref{sec:numerical}) the sufficient condition
$\|J_0^{-1}\|\,\|J_h-J_0\|\le\tfrac12$ would require $h\le2\times10^{-5}$ in the Euclidean norm on
the stage layout and $h\le8\times10^{-4}$ in the endpoint-linearized norm $\|J_0^{-1}\,\cdot\,\|$,
whereas the conclusion $\|J_h^{-1}\|\le2\|J_0^{-1}\|$ holds on every grid used with ratio at most
$1.55$. The gap comes from the curvature of the friction law: the largest entries of $(J_h-J_0)/h$
sit in the Newton--Euler rows against the stage velocities, where the second derivative of the
Brown--McPhee law with $v_s=0.5$ enters, and the same mechanism without joint friction has $C_J$
smaller by a factor of about fifteen and $h_0\approx1.4\times10^{-2}$. The friction law, not the
collocation structure, sets the small-step threshold; Section~\ref{sec:numerical} shows the same
dependence in the observed orders when $v_s$ is decreased.

\subsection{Endpoint defect, closure and inexact Newton}

\begin{lemma}[Gauss endpoint defect]
\label{lem:gauss-endpoint-defect}
Under (H1) there is a constant $C_G$, uniform on the compact tube, such that
\begin{equation}
  \|\Psi_h^G(y)-\varphi_h(y)\|\le C_Gh^7 .
  \label{eq:gauss-endpoint-defect}
\end{equation}
\end{lemma}

\begin{proof}
Let $u$ be the collocation polynomial of Eq.~\eqref{eq:reduced-gauss-stages} and $g$ any
component of the lifted state along $u$, so that $g(t)=\rho(u(t))$ for a smooth chart function
$\rho$ and $g'(t_k)$ equals the lifted stage velocity at $Y_k$; for the rotational component, $g$
is the Lie-algebra increment $\eta(t)$ with $\eta'=J_r^{-1}(\eta)\omega$ along the lifted
trajectory. The reconstruction Eqs.~\eqref{eq:end-trans}--\eqref{eq:end-rot} replaces
$g(t_n+h)-g(t_n)=\int_{t_n}^{t_n+h}g'(t)\,dt$ by the Gauss quadrature $h\sum_ib_ig'(t_i)$. Because
the Gauss weights integrate polynomials of degree at most five exactly, a Peano-kernel argument
gives
\begin{equation}
  \Bigl|\int_{t_n}^{t_n+h}g'(t)\,dt-h\sum_{i=1}^3b_ig'(t_i)\Bigr|
  \le\frac{h^7}{60}\max_{[t_n,t_n+h]}|g^{(7)}| ,
  \label{eq:gauss-quadrature-defect}
\end{equation}
with $g^{(7)}$ bounded uniformly on the tube by (H1) and the smoothness of $u$. Hence each
reconstructed endpoint component equals the lifted endpoint of $u$ up to $O(h^7)$, and $\mathrm{Exp}$
is smooth. Combining with Eq.~\eqref{eq:gauss-local-truncation} and the chart norm equivalence gives
Eq.~\eqref{eq:gauss-endpoint-defect}. The bound Eq.~\eqref{eq:gauss-quadrature-defect} is
machine-checked.
\end{proof}

\begin{lemma}[Endpoint-closure perturbation]
\label{lem:endpoint-closure}
Let $E(z)=0$ be the velocity-level KKT closure subsystem solved by $\mathcal C_h$, and let
$z_u=\Psi_h^G$ be the unclosed endpoint. Under (H2) the right-inverse Newton iteration
$z_{l+1}=z_l-B_\ast E(z_l)$, $z_0=z_u$, converges to a closed endpoint $z_E=\mathcal C_h(z_u)$
with $E(z_E)=0$, and any closed endpoint of the form $z_u+B_\ast w$ in $B_{r_E}(z_\ast)$ satisfies
\begin{equation}
  \|z_E-z_u\|\le 2M_{E}\|E(z_u)\|\le C_Eh^7,\qquad C_E=2M_{E}C_{E,\mathrm{raw}} .
  \label{eq:endpoint-correction-bound}
\end{equation}
\end{lemma}

\begin{proof}
Write $G(w)=E(z_u+B_\ast w)$. Since $DE(z_\ast)B_\ast=I$ and
$\|DE(z)-DE(z_\ast)\|\le L_E\|z-z_\ast\|\le L_Er_E$ on the ball, the mean value inequality gives
$\|G(w_1)-G(w_2)-(w_1-w_2)\|\le M_{E}L_Er_E\|w_1-w_2\|\le\tfrac12\|w_1-w_2\|$ whenever both
arguments stay in the ball. The map $w\mapsto w-G(w)$ is therefore a $\tfrac12$-contraction of the
ball of radius $2\|E(z_u)\|$ into itself, which stays inside $B_{r_E}(z_\ast)$ for $h\le h_0$; its
fixed point $w_\ast$ gives $z_E=z_u+B_\ast w_\ast$ with $E(z_E)=0$. For any root of this form, the
same inequality with $w_2=0$ yields $\|w\|\le2\|E(z_u)\|$, hence
$\|z_E-z_u\|=\|B_\ast w\|\le2M_{E}\|E(z_u)\|$. The raw-defect bound of (H2) completes
Eq.~\eqref{eq:endpoint-correction-bound}.
\end{proof}

\begin{lemma}[Inexact Newton]
\label{lem:inexact-newton}
Under (H2) and (H3),
\begin{equation}
\label{eq:inexact-newton-stage-output-bounds}
\begin{aligned}
  \|\widetilde Z-Z_G\|&\le 2M\eta_h,\\
  \|\mathcal P_h(\widetilde Z)-\mathcal P_h(Z_G)\|&\le C_N\eta_h\le C_Nc_\eta h^7,\qquad C_N=2MM_N .
\end{aligned}
\end{equation}
\end{lemma}

\begin{proof}
With Eq.~\eqref{eq:exact-stage-identity} and the mean value inequality of
Lemma~\ref{lem:stage-residual-defect},
$\|\widetilde Z-Z_G\|\le M\bigl(\|F_h(\widetilde Z)\|+\tfrac{L\rho}{2}\|\widetilde Z-Z_G\|\bigr)$,
so $\|\widetilde Z-Z_G\|\le 2M\eta_h$. The output bound follows from $\|D_Z\mathcal P_h\|\le M_N$ on
the ball and the envelope $\eta_h\le c_\eta h^7$.
\end{proof}

\begin{lemma}[Newton reaches the solver envelope]
\label{lem:newton-envelope}
Under (H1) and (H2) there is $h_0>0$ such that for $0<h\le h_0$ the lifted endpoint predictor
$Z^{(0)}_\ast=(\mathcal L(x_n),\mathcal L(x_n),\mathcal L(x_n))$ lies in $B_{\rho/2}(Z_G)$; for
every predictor $Z^{(0)}\in B_{\rho/2}(Z_G)$ the simplified Newton iteration
$Z^{(k+1)}=Z^{(k)}-J_h^{-1}F_h(Z^{(k)})$ stays in $B_\rho(Z_G)$ and
\begin{equation}
  \|Z^{(k)}-Z_G\|\le 2^{-k}\|Z^{(0)}-Z_G\|,\qquad
  \|F_h(Z^{(k)})\|\le L_F\,2^{-k}\|Z^{(0)}-Z_G\| ,
  \label{eq:newton-envelope-decay}
\end{equation}
where $L_F$ is a Lipschitz constant of $F_h$ on the ball; the same decay holds for the full Newton
iteration of Eq.~\eqref{eq:newton}. Consequently the stopping rule $\|F_h(Z^{(k)})\|\le c_\eta h^7$
is met after $k=O(\log(1/h))$ iterations, and (H3) holds by construction whenever the stage solve
starts in $B_{\rho/2}(Z_G)$.
\end{lemma}

\begin{proof}
The stages of $Z_G$ are $\mathcal L(Y_k)$ with $\|Y_k-y_n\|=O(h)$ on the compact tube, so
$\|Z^{(0)}_\ast-Z_G\|=O(h)$ and $Z^{(0)}_\ast\in B_{\rho/2}(Z_G)$ for small $h$. The linearization
inequality of (H2) makes the simplified Newton map a $\tfrac12$-contraction of $B_\rho(Z_G)$ into
itself with fixed point $Z_G$ (Lemma~\ref{lem:stage-residual-defect}), which gives the first bound
in Eq.~\eqref{eq:newton-envelope-decay}; the second follows from $F_h(Z_G)=0$ and the Lipschitz
bound. For the full Newton iteration, (H2) gives $\|D_ZF_h(Z)-J_h\|\le L\|Z-Z_G\|\le L\rho$ on the
ball, so the perturbation bound of Lemma~\ref{lem:p2-from-p1} yields $\|D_ZF_h(Z)^{-1}\|\le 2M$,
and the standard estimate $\|Z^{(k+1)}-Z_G\|\le ML\|Z^{(k)}-Z_G\|^2\le\tfrac12\|Z^{(k)}-Z_G\|$
gives the same decay. The simplified-Newton decay estimate is machine-checked.
\end{proof}

The predictor of the implementation (Section~\ref{sec:method}) is not $Z^{(0)}_\ast$: it sets the
angular-velocity and angular-acceleration guesses to zero, so its distance to $Z_G$ is of the
order of $\|\omega_n\|+\|\alpha_n\|$ rather than $O(h)$ (about $50$ on the chain,
Table~\ref{tab:predictor-check}), and the first full Newton step from it expands the distance by
$12$ to $18\%$ before the quadratic phase sets in. The lemma therefore covers the implemented run
through the ball membership of (H3), which is confirmed a posteriori on every grid used: the
converged roots are the Gauss roots to roundoff. Replacing the zero angular guesses by the endpoint
values would make the implemented predictor an $O(h)$ predictor and bring the run under the lemma
directly.

The experiments of Section~\ref{sec:numerical} use the fixed tolerance $\tau_N=10^{-11}$ rather
than the $h^7$-scaled rule. On the chain sweep the residual of the accepted iterate is between
$10^{-14}$ and $8\times10^{-12}$ at every step and step size (Table~\ref{tab:e7}); relative to
$h^7$ this is $c_\eta\le10^{-3}$ for $h\ge0.05$ but $c_\eta\approx10^{6}$ at $h=0.003125$, so
the fixed-tolerance runs satisfy (H3) only with a large constant on the finest grids, and the
inexact-Newton term $C_Nc_\eta h^7$ of the theorem is then a pessimistic bound: the observed errors
keep decreasing as $h^6$ down to $10^{-13}$, three orders of magnitude below $2M\|F_h(\widetilde Z)\|$.
One additional Newton iteration per step would bring the accepted residual to roundoff at every
step size; the $h^7$-scaled rule with $c_\eta=10^4$ has been exercised on the chain and gives the
same observed orders as the fixed tolerance.

\subsection{Local-to-global transfer and main theorem}

\begin{lemma}[Local-to-global transfer]
\label{lem:local-global-transfer}
Let $\Psi_h$ be a one-step map on the reduced chart, $t_n=nh$, and $y_{n+1}=\Psi_h(y_n)$ with $y_0$
the exact initial state. Suppose that for $0\le n<N$ with $Nh\le T$ the exact states
$Y_n=\varphi_{t_n}(y_0)$ satisfy $\|\Psi_h(Y_n)-Y_{n+1}\|\le C_\ell h^{p+1}$, that
$\|\Psi_h(y_n)-\Psi_h(Y_n)\|\le(1+C_sh)\|y_n-Y_n\|$ whenever $y_n\in K$, and that
$C_\ell\Gamma_s(T)h^p\le d_K$, where $\Gamma_s(T)=(e^{C_sT}-1)/C_s$ for $C_s>0$ and
$\Gamma_s(T)=T$ for $C_s=0$. Then $y_n\in K$ for all $n\le N$ and
\begin{equation}
\label{eq:local-global-lemma-grid-bound}
  \max_{0\le n\le N}\|y_n-Y_n\|\le C_\ell\Gamma_s(T)h^p .
\end{equation}
If moreover the reporting map $\mathcal R$ (chart to reported position and velocity) is
$C_{\mathcal R}$-Lipschitz on $K$, then
$\max_n\|\mathcal R(y_n)-\mathcal R(Y_n)\|\le C_{\mathcal R}C_\ell\Gamma_s(T)h^p$.
\end{lemma}

\begin{proof}
Let $e_n=\|y_n-Y_n\|$. By induction, as long as $y_n\in K$,
$e_{n+1}\le(1+C_sh)e_n+C_\ell h^{p+1}$, hence $e_n\le C_\ell h^{p+1}\sum_{k<n}(1+C_sh)^k$. For
$C_s>0$ the geometric sum equals $((1+C_sh)^n-1)/(C_sh)\le(e^{C_snh}-1)/(C_sh)\le\Gamma_s(T)/h$,
and for $C_s=0$ it equals $n\le T/h$. Thus $e_n\le C_\ell\Gamma_s(T)h^p\le d_K$, which keeps $y_n$
in the tube and closes the induction. The reporting bound is the Lipschitz property. The lemma is
machine-checked.
\end{proof}

\begin{definition}[Branch-selected one-step map]
\label{def:accepted-branch-map}
The one-step map $\Psi_h$ is defined on the set $\mathcal D_h\subset K$ of reduced states for which
the transition satisfies (H2) and (H3): the computed iterate lies in $B_\rho(Z_G)$ with residual at
most $\eta_h$, the endpoint is reconstructed by Eqs.~\eqref{eq:end-trans}--\eqref{eq:end-rot} and
closed by $\mathcal C_h$, and $\Psi_h(y)=\mathcal P_h(\widetilde Z)$.
\end{definition}

\begin{theorem}[Sixth-order error bound]
\label{thm:g6fullva-order}
Let Assumption~\ref{ass:regularity} hold and let $C_{\rm loc}=C_G+C_E+C_Nc_\eta$ with $C_G$, $C_E$,
$C_N$ from Lemmas~\ref{lem:gauss-endpoint-defect}, \ref{lem:endpoint-closure}
and~\ref{lem:inexact-newton}. Then there is $h_0>0$ such that for $0<h\le h_0$ and every
$y\in\mathcal D_h$,
\begin{equation}
\label{eq:g6fva-local-defect-theorem}
  \|\Psi_h(y)-\varphi_h(y)\|\le C_{\rm loc}h^7 .
\end{equation}
Consequently, for every grid $t_n=nh$, $0\le n\le N_h$, $N_hh\le T$, whose transitions lie in
$\mathcal D_h$ and whose numerical trajectory $y_{n+1}=\Psi_h(y_n)$ starts from the exact reduced
initial state,
\begin{equation}
\label{eq:g6fva-reduced-grid-bound}
  \max_{0\le n\le N_h}\|y_n-y(t_n)\|\le C_{\rm red}h^6,\qquad C_{\rm red}=C_{\rm loc}\Gamma_s(T),
\end{equation}
and the reported position--velocity map $\mathcal R(y)=(q(y),v(y))$ satisfies
\begin{equation}
\label{eq:g6fva-reported-grid-bound}
  \max_{0\le n\le N_h}\|\mathcal R(y_n)-\mathcal R(y(t_n))\|\le C_{qv}h^6,\qquad
  C_{qv}=C_{\mathcal R}C_{\rm red} .
\end{equation}
The constants depend on the tube, on $T$ and on the constants of Assumption~\ref{ass:regularity},
including $c_\eta$, but not on $h$, $N_h$ or the linear-algebra backend. Under (H1) alone the same
conclusions hold for $h\le h_0$ whenever the stage solve is started in $B_{\rho/2}(Z_G)$ and
stopped by the rule $\|F_h\|\le c_\eta h^7$ (Lemmas~\ref{lem:p2-from-p1}
and~\ref{lem:newton-envelope}).
\end{theorem}

\begin{proof}
Fix $y\in\mathcal D_h$ and write $\Psi_h^G=\mathcal E_h(Z_G)$, $\Psi_h^E=\mathcal P_h(Z_G)$ and
$\Psi_h(y)=\mathcal P_h(\widetilde Z)$. By Lemma~\ref{lem:stage-residual-defect} the exact stage
root on the branch is $Z_G$, so there is no stage-root perturbation term, and
\begin{equation}
\label{eq:three-term-local-defect-sum}
\begin{aligned}
  \|\Psi_h(y)-\varphi_h(y)\|
  &\le\|\Psi_h^G(y)-\varphi_h(y)\|+\|\Psi_h^E(y)-\Psi_h^G(y)\|\\
  &\quad+\|\Psi_h(y)-\Psi_h^E(y)\|
  \le C_Gh^7+C_Eh^7+C_Nc_\eta h^7
\end{aligned}
\end{equation}
by Lemmas~\ref{lem:gauss-endpoint-defect}, \ref{lem:endpoint-closure}
and~\ref{lem:inexact-newton}, after shrinking $h_0$ so that all three apply on the tube. This is
Eq.~\eqref{eq:g6fva-local-defect-theorem}. Choose $h_0$ additionally so that
$C_{\rm loc}\Gamma_s(T)h_0^6\le d_K$. Lemma~\ref{lem:local-global-transfer} with $p=6$,
$C_\ell=C_{\rm loc}$ and the stability scale Eq.~\eqref{eq:assumption-stability-scale} gives
Eq.~\eqref{eq:g6fva-reduced-grid-bound}; the reporting-map clause gives
Eq.~\eqref{eq:g6fva-reported-grid-bound}. The last sentence follows from
Lemmas~\ref{lem:p2-from-p1} and~\ref{lem:newton-envelope}, which supply (H2) and (H3) from (H1) for
small steps.
\end{proof}

Two remarks on the scope of the theorem. First, it is an estimate for the reduced state and for the
reported positions and velocities; it says nothing about the accuracy of the stage multipliers or
reaction forces beyond what follows from the smooth dependence of the lift on the reduced state, and
it does not turn a small constraint residual on a closed-loop mechanism into a trajectory-error
bound, which would need a stability constant for that mechanism. Second, because the stage residual
vanishes identically at the lifted Gauss stage, the order of the method is that of reduced Gauss
collocation; the absolute-coordinate formulation contributes the exact three-level constraint
enforcement, the multipliers and reaction forces at every stage, the friction law through the
multipliers, and the Lie-group endpoint reconstruction, but no additional stage-level
approximation. The same argument gives order $2s$ for the $s$-stage Gauss member with the same
stage system, which Section~\ref{sec:numerical} confirms for $s=1,2$.

\paragraph{Formal-order comparison with time finite elements}
The Lie-group time finite-element family of~\cite{chaturvedi2026tfe} with Gauss--Lobatto
polynomials of degree $m$ has formal order $2m-1$~\cite{kim2017higher}; its $m=3$ member is a
fifth-order formula. The comparison of the present method with that family is at this formal level
only: order six against order five for a stage system of comparable size. We have not implemented
the time finite-element method and make no statement about its accuracy or cost on its own test
problems.

\section{Numerical experiments}
\label{sec:numerical}

All experiments use the implementation of Section~\ref{sec:method} with the fixed Newton tolerance
$\tau_N=10^{-11}$ and the dense automatic-differentiation Jacobian; the scripts that produce every
number and figure of this section are listed in \ref{app:repro}. Errors are measured
against a reference solution computed with the same method at a step size $h_{\rm ref}$ much
smaller than the finest step of the sweep; the difference between the references at $h_{\rm ref}$
and $2h_{\rm ref}$ is reported as the resolution floor of the reference, and a fitted order uses only
the step sizes whose errors exceed that floor by a factor of $100$. Position and velocity errors are
Frobenius norms over the bodies at the final time; the orientation error is the largest geodesic
angle between the numerical and reference rotations, and the angular-velocity error the Frobenius
norm of the body-frame difference.

\subsection{Convergence on the frictional chain}
\label{sec:e1}

The chain of Table~\ref{tab:chain} with $v_s=0.5$ is integrated over its regular branch $[0,0.5]$
with $h=0.1/2^k$, $k=0,\dots,5$, against the reference at $h_{\rm ref}=0.1/256$ ($1280$ steps).
Table~\ref{tab:e1} and Figure~\ref{fig:e1} show the result. All four error measures decrease by
$5.5$ to $6$ orders of magnitude over the sweep; the fitted orders over the five step pairs above
the floor are $6.25$ (position), $6.28$ (orientation), $6.11$ (velocity) and $6.27$ (angular
velocity), and the consecutive pairwise orders settle at $5.98$ and $5.99$ on the two finest
pairs. The endpoint velocity-level constraint norm is at roundoff on every step because of the
closure; the position-level norm, which is not projected, decreases as $h^6$ from $10^{-7}$ to
$3\times10^{-15}$ as the theory predicts. The Newton solve takes $6.0$ iterations per step on the
coarse grids and $5.7$ on the fine ones.

\begin{table}[!htb]
\centering
\caption{Convergence on the frictional chain, $T=0.5$, $v_s=0.5$: errors at $t=T$ against the
$h_{\rm ref}=0.1/256$ reference, pairwise order of the position error to the next finer step, and
maximal endpoint constraint norms at position ($\Phi$) and velocity ($\Phi_qv$) level; the number of steps is $0.5/h$. Resolution
floor of the reference: $3\times10^{-15}$ (position), $3\times10^{-15}$ (orientation),
$2\times10^{-14}$ (velocity), $1\times10^{-14}$ (angular velocity).}
\label{tab:e1}
\scriptsize
\setlength{\tabcolsep}{2.5pt}
\begin{tabular}{@{}lrllllrll@{}}
\toprule
$h$ & Newton/step & position & orientation & velocity & ang.\ velocity & order & $\|\Phi\|$ & $\|\Phi_qv\|$ \\
\midrule
$0.1$      & 6.00 & $2.77\mathrm{e}{-4}$  & $8.77\mathrm{e}{-5}$  & $4.96\mathrm{e}{-4}$  & $1.10\mathrm{e}{-3}$  & 6.95 & $9.8\mathrm{e}{-8}$  & $7.7\mathrm{e}{-16}$ \\
$0.05$     & 6.00 & $2.24\mathrm{e}{-6}$  & $7.98\mathrm{e}{-7}$  & $2.18\mathrm{e}{-6}$  & $9.47\mathrm{e}{-6}$  & 4.58 & $1.5\mathrm{e}{-9}$  & $1.3\mathrm{e}{-15}$ \\
$0.025$    & 5.90 & $9.38\mathrm{e}{-8}$  & $2.82\mathrm{e}{-8}$  & $1.95\mathrm{e}{-7}$  & $3.59\mathrm{e}{-7}$  & 7.64 & $2.4\mathrm{e}{-11}$ & $1.2\mathrm{e}{-15}$ \\
$0.0125$   & 5.72 & $4.69\mathrm{e}{-10}$ & $1.43\mathrm{e}{-10}$ & $9.42\mathrm{e}{-10}$ & $1.82\mathrm{e}{-9}$  & 5.98 & $3.7\mathrm{e}{-13}$ & $9.8\mathrm{e}{-14}$ \\
$0.00625$  & 5.72 & $7.45\mathrm{e}{-12}$ & $2.28\mathrm{e}{-12}$ & $1.49\mathrm{e}{-11}$ & $2.89\mathrm{e}{-11}$ & 5.99 & $6.3\mathrm{e}{-15}$ & $9.6\mathrm{e}{-14}$ \\
$0.003125$ & 5.72 & $1.17\mathrm{e}{-13}$ & $3.55\mathrm{e}{-14}$ & $2.33\mathrm{e}{-13}$ & $4.44\mathrm{e}{-13}$ & --   & $3.0\mathrm{e}{-15}$ & $8.2\mathrm{e}{-15}$ \\
\midrule
fitted order & & 6.25 & 6.28 & 6.11 & 6.27 & & & \\
\bottomrule
\end{tabular}
\end{table}

\paperfig{\figpath/Figure_3_e1_convergence.png}
{Convergence on the frictional chain over the regular branch $[0,0.5]$: errors at $t=T$ (left) and
over all grid points (right) against the $h_{\rm ref}=0.1/256$ reference, with the slope-6 line
anchored at the finest position error and the resolution floor of the reference.}
{fig:e1}

\paragraph{Exact stage identity and constants of the analysis}
Three diagnostics of Section~\ref{sec:order} are evaluated on the chain runs.
(i)~Each converged stage vector, obtained with a Newton tolerance of $10^{-13}$, is read in the
joint-coordinate chart and the reduced collocation defects $\Delta_i[s_j]$, $\Delta_i[\dot s_j]$,
$\Delta_i[\theta_j]$, $\Delta_i[\dot\theta_j]$ of Eq.~\eqref{eq:collocation-defect-operator} are
evaluated together with the component of $d_j,\dot d_j,\ddot d_j$ orthogonal to $n_j$: all are at
roundoff (Table~\ref{tab:exact-identity-check}), as Lemma~\ref{lem:exact-stage-identity} states.
(ii)~The constants of Lemma~\ref{lem:p2-from-p1} are measured from the implemented Jacobian at the
converged stage and at the root of the $h=0$ stage system (Table~\ref{tab:p2-constants-check}): the
conclusion $\|J_h^{-1}\|\le2\|J_0^{-1}\|$ holds with ratio at most $1.55$, while the sufficient
Neumann condition would need $h\le2\times10^{-5}$ in the Euclidean norm and $h\le8\times10^{-4}$ in
the endpoint-linearized norm; without joint friction the threshold rises to $1.4\times10^{-2}$.
(iii)~The distances of the two predictors to $Z_G$ (Table~\ref{tab:predictor-check}) confirm the
discussion after Lemma~\ref{lem:newton-envelope}: the lifted endpoint predictor is $O(h)$ from the
Gauss stage, the implemented predictor is at distance about $50$ and its first Newton step expands
that distance by up to $18\%$.

\begin{table}[!htb]
\centering
\caption{Exact stage identity at the converged stages of the chain (Newton tolerance $10^{-13}$,
window $[0,0.08]$): maximal implemented residual, maximal reduced collocation defect and maximal
component of $d_j,\dot d_j,\ddot d_j$ orthogonal to $n_j$.}
\label{tab:exact-identity-check}
\small
\begin{tabular}{P{0.07\linewidth}P{0.08\linewidth}P{0.13\linewidth}P{0.17\linewidth}P{0.17\linewidth}P{0.17\linewidth}}
\toprule
$h$ & steps & Newton total & max stage residual & max reduced defect & max $\perp n_j$ component \\
\midrule
0.04 & 2 & 12 & $6.8\times10^{-14}$ & $9.0\times10^{-15}$ & $1.5\times10^{-15}$ \\
0.02 & 4 & 24 & $1.5\times10^{-14}$ & $5.5\times10^{-16}$ & $1.9\times10^{-15}$ \\
0.01 & 8 & 48 & $1.6\times10^{-14}$ & $5.8\times10^{-16}$ & $2.0\times10^{-15}$ \\
\bottomrule
\end{tabular}
\end{table}

\begin{table}[!htb]
\centering
\caption{Constants of Lemma~\ref{lem:p2-from-p1} on the chain: $M_0=\max_n\|J_0^{-1}\|$,
$C_J=\max_n\|J_h-J_0\|/h$, observed inverse ratio $\|J_h^{-1}\|/\|J_0^{-1}\|$ and Neumann threshold
$h_0=1/(2M_0C_J)$, in the Euclidean norm on the stage layout (E) and in the endpoint-linearized
norm $\|J_0^{-1}\,\cdot\,\|$ (L, where $M_0=1$); the last two columns repeat the L-norm quantities
for the same mechanism without joint friction.}
\label{tab:p2-constants-check}
\scriptsize
\setlength{\tabcolsep}{3pt}
\begin{tabular}{@{}lrlrllllr@{}}
\toprule
$h$ & $M_0$ (E) & $C_J$ (E) & ratio & $h_0$ (E) & $C_J$ (L) & $h_0$ (L) & $C_J$ (L, no fr.) & $h_0$ (L, no fr.) \\
\midrule
0.04 & 18.2 & $7.5\mathrm{e}{2}$ & 1.55 & $3.7\mathrm{e}{-5}$ & $5.1\mathrm{e}{2}$ & $9.8\mathrm{e}{-4}$ & 34.1 & $1.5\mathrm{e}{-2}$ \\
0.02 & 24.3 & $8.2\mathrm{e}{2}$ & 1.26 & $2.5\mathrm{e}{-5}$ & $5.6\mathrm{e}{2}$ & $8.9\mathrm{e}{-4}$ & 35.5 & $1.4\mathrm{e}{-2}$ \\
0.01 & 28.4 & $9.6\mathrm{e}{2}$ & 1.18 & $1.8\mathrm{e}{-5}$ & $6.5\mathrm{e}{2}$ & $7.6\mathrm{e}{-4}$ & 36.5 & $1.4\mathrm{e}{-2}$ \\
\bottomrule
\end{tabular}
\end{table}

\begin{table}[!htb]
\centering
\caption{Predictor distances for Lemma~\ref{lem:newton-envelope} on the chain (Euclidean norm on
the stage layout): distance of the lifted endpoint predictor $Z^{(0)}_\ast$ and of the implemented
predictor to $Z_G$, and the ratio $\|Z^{(1)}-Z_G\|/\|Z^{(0)}-Z_G\|$ of the first full Newton step
from the implemented predictor.}
\label{tab:predictor-check}
\small
\begin{tabular}{P{0.08\linewidth}P{0.20\linewidth}P{0.22\linewidth}P{0.20\linewidth}}
\toprule
$h$ & lifted predictor & implemented predictor & first-step ratio \\
\midrule
0.04 & 3.23 & 50.2 & 1.12 \\
0.02 & 2.11 & 50.3 & 1.16 \\
0.01 & 1.07 & 50.4 & 1.18 \\
\bottomrule
\end{tabular}
\end{table}

\paragraph{Solver envelope}
Table~\ref{tab:e7} records, for every grid of the sweep, the largest residual of the accepted Newton
iterate under the production stopping rule and its ratio to $h^7$, which is the constant $c_\eta$
with which that grid satisfies (H3). The accepted residual is set by the tolerance, not by $h$, so
the ratio grows from $10^{-4}$ to $10^{6}$ across the sweep; as discussed after
Lemma~\ref{lem:newton-envelope}, the resulting bound $2M\|F_h(\widetilde Z)\|$ on the stage
perturbation is pessimistic by three orders of magnitude relative to the observed errors. The
endpoint position-level constraint norm of Table~\ref{tab:e1} is the quantity that actually
dominates the deviation of the numerical trajectory from the constraint manifold; it decreases as
$h^6$ as the theorem requires.

\begin{table}[!htb]
\centering
\caption{Solver envelope on the chain sweep ($\tau_N=10^{-11}$): largest residual norm of the
accepted iterate over all steps, its ratio to $h^7$ (the $c_\eta$ of (H3) for that grid), and the
largest components of the stage displacements $d_j,\dot d_j,\ddot d_j$ orthogonal to $n_j$.}
\label{tab:e7}
\footnotesize
\begin{tabular}{@{}lrlll@{}}
\toprule
$h$ & Newton/step & $\max_n\|F_h(\widetilde Z_n)\|$ & ratio to $h^7$ ($c_\eta$) & max $\perp n_j$ component \\
\midrule
$0.1$      & 6.00 & $8.1\mathrm{e}{-12}$ & $8.1\mathrm{e}{-5}$ & $5.6\mathrm{e}{-13}$ \\
$0.05$     & 6.00 & $3.5\mathrm{e}{-13}$ & $4.5\mathrm{e}{-4}$ & $3.2\mathrm{e}{-15}$ \\
$0.025$    & 5.90 & $8.8\mathrm{e}{-12}$ & $1.4$               & $6.4\mathrm{e}{-14}$ \\
$0.0125$   & 5.72 & $2.7\mathrm{e}{-12}$ & $5.7\mathrm{e}{1}$  & $4.3\mathrm{e}{-15}$ \\
$0.00625$  & 5.72 & $3.1\mathrm{e}{-12}$ & $8.3\mathrm{e}{3}$  & $4.1\mathrm{e}{-15}$ \\
$0.003125$ & 5.72 & $3.4\mathrm{e}{-12}$ & $1.2\mathrm{e}{6}$  & $5.3\mathrm{e}{-15}$ \\
\bottomrule
\end{tabular}
\end{table}

\subsection{The Gauss family: orders two, four and six at the same Newton work}
\label{sec:e2}

Replacing the three-stage tableau by the two-stage Gauss tableau (order four) or by the implicit
midpoint rule (order two) leaves the stage system of Eq.~\eqref{eq:g6fullva-expanded-stage}
unchanged except for its size ($88$ and $44$ unknowns). Figure~\ref{fig:e2} and
Table~\ref{tab:e2} compare the three members on the chain over $[0,0.5]$ against the same
reference. The fitted position orders are $2.10$, $4.17$ and $6.25$, in agreement with the remark
after Theorem~\ref{thm:g6fullva-order}. The Newton iteration count per step is the same for the
three members ($5.7$ to $6.0$), so at a given step size the work differs only by the size of the
dense factorization; in the present Python implementation the wall time per step differs by at most
a third across the family, being dominated by interpreter overhead. At $h=0.0125$ the three-stage
member is $1.6\times10^{5}$ times more accurate than the one-stage and $310$ times more accurate
than the two-stage member, and at $h=0.003125$ the factors are $4\times10^{7}$ and
$5\times10^{3}$; to reach the accuracy of the three-stage member at $h=0.05$ ($2\times10^{-6}$)
the one-stage member needs $h\approx2\times10^{-3}$, that is, about $23$ times more steps and
Newton iterations.

\begin{table}[!htb]
\centering
\caption{Gauss family with the same stage system on the chain, $T=0.5$: position and velocity
errors at $t=T$ and total Newton iterations. The wall time of the first row of each member includes
just-in-time compilation and is omitted.}
\label{tab:e2}
\scriptsize
\setlength{\tabcolsep}{3pt}
\begin{tabular}{@{}lllrlllr@{}}
\toprule
$h$ & \multicolumn{3}{l}{1 stage (order 2)} & \multicolumn{3}{l}{2 stages (order 4)} & \\
    & position & velocity & Newton & position & velocity & Newton & \\
\midrule
$0.1$      & $1.16\mathrm{e}{-2}$ & $6.53\mathrm{e}{-2}$ & 31   & $1.48\mathrm{e}{-3}$ & $3.23\mathrm{e}{-3}$ & 30 & \\
$0.05$     & $1.16\mathrm{e}{-3}$ & $1.21\mathrm{e}{-2}$ & 60   & $6.23\mathrm{e}{-5}$ & $8.38\mathrm{e}{-5}$ & 60 & \\
$0.025$    & $2.95\mathrm{e}{-4}$ & $2.97\mathrm{e}{-3}$ & 115  & $1.61\mathrm{e}{-6}$ & $2.81\mathrm{e}{-6}$ & 117 & \\
$0.0125$   & $7.41\mathrm{e}{-5}$ & $7.39\mathrm{e}{-4}$ & 229  & $1.45\mathrm{e}{-7}$ & $2.66\mathrm{e}{-7}$ & 229 & \\
$0.00625$  & $1.87\mathrm{e}{-5}$ & $1.85\mathrm{e}{-4}$ & 457  & $9.09\mathrm{e}{-9}$ & $1.67\mathrm{e}{-8}$ & 458 & \\
$0.003125$ & $4.68\mathrm{e}{-6}$ & $4.61\mathrm{e}{-5}$ & 915  & $5.69\mathrm{e}{-10}$ & $1.04\mathrm{e}{-9}$ & 915 & \\
$0.0015625$ & $1.17\mathrm{e}{-6}$ & $1.15\mathrm{e}{-5}$ & 1829 & $3.55\mathrm{e}{-11}$ & $6.51\mathrm{e}{-11}$ & 1830 & \\
$0.00078125$ & $2.93\mathrm{e}{-7}$ & $2.88\mathrm{e}{-6}$ & 3658 & -- & -- & -- & \\
\midrule
fitted order & 2.10 & 2.04 & & 4.17 & 4.17 & & \\
\bottomrule
\end{tabular}
\end{table}

\paperfig{\figpath/Figure_4_e2_gauss_family.png}
{Gauss family with the same FullVA stage system on the chain, $T=0.5$: position error against step
size (left) and against total Newton iterations (right). Three stages (order six) are the method of
this paper.}
{fig:e2}

\subsection{Sharpness of the friction law}
\label{sec:e3}

The Brown--McPhee law is smooth for every $v_s>0$, but its derivatives grow like $v_s^{-k}$, and
Lemma~\ref{lem:p2-from-p1} predicts a small-step threshold that shrinks accordingly.
Figure~\ref{fig:e3} and Table~\ref{tab:e3} repeat the convergence sweep of Section~\ref{sec:e1}
for $v_s\in\{0.5,0.2,0.1,0.05,0.02\}$, each against its own reference at $h_{\rm ref}=0.1/256$.
Two things happen as $v_s$ decreases. The fitted order over the whole sweep drops from $6.25$ to
$2.47$, because the coarse grids are in a pre-asymptotic regime in which the error is dominated by
the friction transition; and the step size at which the pairwise order returns to six moves from
$h=0.1$ to $h\approx3\times10^{-3}$, where all five cases show pairwise orders between $5.3$ and
$6.7$ on the finest pair. The measured Neumann threshold $h_0$ decreases from $7.6\times10^{-4}$ to
$7.1\times10^{-6}$ over the same range, so the analysis predicts the trend correctly but its
sufficient condition is two to three orders of magnitude more conservative than the observed
recovery step. Every run in the sweep converged, with $5.7$ to $6.4$ Newton iterations per step.

\begin{table}[!htb]
\centering
\caption{Friction sweep on the chain, $T=0.5$: fitted position and velocity orders over the six
step sizes, pairwise position orders from $h=0.1$ to $h=0.003125$, coarsest step at which the
pairwise order reaches $5.5$, and the Neumann threshold $h_0=1/(2C_J)$ in the endpoint-linearized
norm measured over the first ten steps at $h=0.01$.}
\label{tab:e3}
\scriptsize
\setlength{\tabcolsep}{3.5pt}
\begin{tabular}{@{}lrrlll@{}}
\toprule
$v_s$ & position order & velocity order & pairwise position orders & recovery $h$ & $h_0$ (L) \\
\midrule
0.5  & 6.25 & 6.11 & 6.95, 4.58, 7.64, 5.98, 5.99 & 0.1     & $7.6\mathrm{e}{-4}$ \\
0.2  & 4.84 & 4.74 & 2.11, 4.54, 2.74, 8.57, 5.85 & 0.0125  & $1.4\mathrm{e}{-4}$ \\
0.1  & 4.74 & 4.78 & 0.47, 5.71, 4.14, 5.88, 6.72 & 0.05    & $6.7\mathrm{e}{-5}$ \\
0.05 & 3.66 & 3.66 & 0.17, 4.11, 4.47, 3.46, 5.27 & --      & $2.1\mathrm{e}{-5}$ \\
0.02 & 2.47 & 2.48 & $-0.22$, 2.29, 2.90, 1.55, 6.20 & 0.00625 & $7.1\mathrm{e}{-6}$ \\
\bottomrule
\end{tabular}
\end{table}

\paperfig{\figpath/Figure_5_e3_friction_sweep.png}
{Friction sweep on the chain: position error against step size for five Stribeck velocities (left)
and fitted orders against $v_s$ (right).}
{fig:e3}

\subsection{The four benchmark mechanisms}
\label{sec:e4}

The public benchmark suite of~\cite{kissel2022absolute} consists of a driven single pendulum, a
free double pendulum released from the horizontal, and two driven closed loops, a four-link
mechanism and a slider-crank (\ref{app:lower-pair}). Three of the four are kinematically
driven: their motion is fixed by the drive constraint and the dynamics only determines the
reactions. \method{} is run on all four over $T=1$ with $h=0.1/2^k$ (Figure~\ref{fig:e4},
Table~\ref{tab:e4}). For the driven single pendulum, whose exact motion is known in closed form,
the final-state position errors are $5.9\times10^{-10}$, $9.3\times10^{-12}$ and
$1.5\times10^{-13}$ at $h=0.1$, $0.05$ and $0.025$ (pairwise orders $6.0$ and $6.0$) and at
roundoff, $10^{-14}$ to $10^{-15}$, for every finer step; its velocity errors scatter between
$10^{-15}$ and $4\times10^{-9}$ without trend in $h$, at the level set by the Newton tolerance of
that harness, which does not project the endpoint velocities. For the two closed loops the trajectory errors against the
$h=0.1/128$ reference are between $6\times10^{-16}$ and $1.3\times10^{-14}$ at every step size
from $0.05$ down to $0.0016$: the driven motion is reproduced to roundoff, the loop-closure
constraints included, and no order can be measured. The double pendulum is the only genuine
dynamics test of the suite. Its final-state errors against the $h=0.1/128$ reference are
$4.6\times10^{-5}$, $1.6\times10^{-7}$, $1.2\times10^{-9}$, $1.9\times10^{-11}$ and
$3.1\times10^{-13}$ for $h=0.1$ to $0.00625$, with pairwise orders $8.2$, $7.0$, $6.0$, $6.0$ in
position and $9.7$, $7.0$, $6.0$, $5.9$ in velocity; the coarsest pair is pre-asymptotic, the rest
is the sixth order of Theorem~\ref{thm:g6fullva-order}. At $h=0.003125$ the error,
$2.3\times10^{-14}$, has reached the resolution floor of the reference ($2\times10^{-14}$,
the difference between the $h=0.1/128$ and $h=0.1/64$ runs).

\paperfig{\figpath/Figure_6_mechanisms.png}
{The four benchmark mechanisms of the public suite~\cite{kissel2022absolute}; joint data in
\ref{app:lower-pair}.}
{fig:mechanisms}

\begin{table}[!htb]
\centering
\caption{The four benchmark mechanisms over $T=1$: final-state position and velocity errors (single
pendulum: against the analytic driven motion; double pendulum: against the $h=0.1/128$ reference)
and trajectory $L^\infty$ errors of the closed loops against their $h=0.1/128$ reference. The
pairwise order refers to the position error and the next finer step; the double-pendulum reference
has a resolution floor of $2\times10^{-14}$.}
\label{tab:e4}
\scriptsize
\setlength{\tabcolsep}{3pt}
\begin{tabular}{@{}lllrllr@{}}
\toprule
$h$ & \multicolumn{3}{l}{single pendulum (driven, analytic reference)} & \multicolumn{3}{l}{double pendulum (free)} \\
    & position & velocity & order & position & velocity & order \\
\midrule
$0.1$      & $5.94\mathrm{e}{-10}$ & $2.43\mathrm{e}{-9}$  & 6.0 & $4.59\mathrm{e}{-5}$  & $1.56\mathrm{e}{-4}$  & 8.2 \\
$0.05$     & $9.32\mathrm{e}{-12}$ & $3.83\mathrm{e}{-11}$ & 6.0 & $1.58\mathrm{e}{-7}$  & $1.85\mathrm{e}{-7}$  & 7.0 \\
$0.025$    & $1.45\mathrm{e}{-13}$ & $9.49\mathrm{e}{-11}$ & --  & $1.24\mathrm{e}{-9}$  & $1.48\mathrm{e}{-9}$  & 6.0 \\
$0.0125$   & $1.08\mathrm{e}{-14}$ & $6.77\mathrm{e}{-10}$ & --  & $1.93\mathrm{e}{-11}$ & $2.29\mathrm{e}{-11}$ & 6.0 \\
$0.00625$  & $1.58\mathrm{e}{-14}$ & $3.73\mathrm{e}{-9}$  & --  & $3.10\mathrm{e}{-13}$ & $3.86\mathrm{e}{-13}$ & 3.8 \\
$0.003125$ & $1.33\mathrm{e}{-15}$ & $1.05\mathrm{e}{-15}$ & --  & $2.29\mathrm{e}{-14}$ & $6.15\mathrm{e}{-14}$ & -- \\
\midrule
 & \multicolumn{3}{l}{four-link (driven loop)} & \multicolumn{3}{l}{slider-crank (driven loop)} \\
 & position & velocity & & position & velocity & \\
\midrule
$0.05$--$0.0016$ & $0.9$--$1.3\mathrm{e}{-14}$ & $1.5$--$4.0\mathrm{e}{-14}$ & -- & $0.6$--$1.2\mathrm{e}{-15}$ & $1.4$--$2.8\mathrm{e}{-15}$ & -- \\
\bottomrule
\end{tabular}
\end{table}

\paperfig{\figpath/Figure_7_e4_benchmarks.png}
{\method{} on the four benchmark mechanisms over $T=1$. The three driven mechanisms are reproduced at
roundoff for every step size; the double pendulum shows the sixth order.}
{fig:e4}

\subsection{Comparison with the public absolute-coordinate codes}
\label{sec:e5}

The public codes of the benchmark suite (the 2021 rA, rp and r$\varepsilon$ formulations of
\cite{kissel2022absolute} and the 2022 half-implicit rA and rA-half codes of
\cite{fang2023halfimplicit}) are run as published, in their dynamics mode with their own tolerance
policies, on the double pendulum over the public horizon $T=3$. Every error in
Table~\ref{tab:e5} is the final-state $L^\infty$ difference to one common reference, the local
\method{} run at $h=10^{-4}$ ($30\,000$ steps). That the public model and the local one are the same
mechanism is checked on the short horizon $T=0.1$, where the public rA final state differs from
the local one by $3.7\times10^{-3}$ at $h=0.1/16$ and $9.2\times10^{-4}$ at $h=0.1/64$, a
first-order decrease and not a constant offset.

Over $T=3$ the double pendulum amplifies perturbations strongly, and the public codes do not
resolve it at any of the step sizes: the rA, r$\varepsilon$ and 2022 rA codes give identical
results with position errors from $1.8$ at $h=0.1$ to $0.17$ at $h=0.001$ (pairwise orders
$0.2$ to $0.9$, approaching first order only below $h=0.002$); the rp code fails to converge for
$h\ge0.025$ and is at first order below; the half-implicit rA-half code does not converge at all
on this horizon (errors between $0.9$ and $5.8$). \method{} reaches $4.3\times10^{-4}$ at
$h=0.1$, $2.1\times10^{-11}$ at $h=0.01$ and the resolution floor of the reference,
$2\times10^{-13}$, at $h\le0.002$, with pairwise orders $7.2$, $9.2$, $6.0$ and $6.1$ on the first
four pairs. The Newton counts are $3.0$ to $4.5$ iterations per step for \method{} against $1.8$
to $6.8$ for the 2021 codes and $3$ to $32$ for the 2022 codes; wall times are not comparable
across the implementations (the local harness is a Python reference implementation with a dense
automatic-differentiation Jacobian) and are omitted.

On the three driven mechanisms the public horizon adds nothing to Section~\ref{sec:e4}:
\method{} reproduces the driven motion to $10^{-14}$ over $T=3$ at $h=10^{-2}$, $10^{-3}$ and
$10^{-4}$, while the public rA code, measured against its own kinematic-mode reference, has
position errors below $10^{-9}$ but velocity errors that decrease at first order,
$1.1\times10^{-1}$, $1.1\times10^{-2}$, $1.1\times10^{-3}$ on the four-link mechanism and
$8.5\times10^{-3}$, $8.7\times10^{-4}$, $8.8\times10^{-5}$ on the slider-crank.

\begin{table}[!htb]
\centering
\caption{Double pendulum over the public horizon $T=3$: final-state $L^\infty$ errors against the
common \method{} reference at $h=10^{-4}$. The 2021 r$\varepsilon$ and the 2022 half-implicit rA
codes return the same numbers as the 2021 rA code. ``fail'': the public Newton iteration did not
converge. The reference has a resolution floor of about $2\times10^{-13}$.}
\label{tab:e5}
\scriptsize
\setlength{\tabcolsep}{3.5pt}
\begin{tabular}{@{}lllllll@{}}
\toprule
$h$ & \multicolumn{2}{l}{\method{}} & \multicolumn{2}{l}{2021 public rA} & 2021 public rp & 2022 rA-half \\
    & position & velocity & position & velocity & position & position \\
\midrule
$0.1$    & $4.29\mathrm{e}{-4}$  & $1.51\mathrm{e}{-3}$  & $1.84$ & $6.04$ & fail   & $5.76$ \\
$0.05$   & $3.03\mathrm{e}{-6}$  & $9.84\mathrm{e}{-6}$  & $1.58$ & $5.03$ & fail   & $5.01$ \\
$0.025$  & $5.34\mathrm{e}{-9}$  & $1.72\mathrm{e}{-8}$  & $1.34$ & $4.13$ & fail   & $1.34$ \\
$0.0125$ & $8.12\mathrm{e}{-11}$ & $2.66\mathrm{e}{-10}$ & $1.06$ & $3.43$ & $0.749$ & $0.903$ \\
$0.01$   & $2.10\mathrm{e}{-11}$ & $7.09\mathrm{e}{-11}$ & $0.958$ & $3.40$ & $1.25$ & $1.12$ \\
$0.002$  & $2.56\mathrm{e}{-13}$ & $1.76\mathrm{e}{-12}$ & $0.319$ & $1.10$ & $0.233$ & $4.26$ \\
$0.001$  & $2.16\mathrm{e}{-13}$ & $1.66\mathrm{e}{-12}$ & $0.169$ & $0.532$ & $0.113$ & $1.92$ \\
\bottomrule
\end{tabular}
\end{table}

\paperfig{\figpath/Figure_8_e5_double_pendulum.png}
{Double pendulum over the public horizon $T=3$: final position (left) and velocity (right) errors
of \method{} and of the public absolute-coordinate codes against the common $h=10^{-4}$ reference.}
{fig:e5}


\subsection{Constraint drift and energy over the regular branch}
\label{sec:e6}

Figure~\ref{fig:e6} follows the chain over the whole regular branch $[0,0.5]$ at $h=0.01$ in two
variants: the frictional benchmark, and a conservative variant with $\mu_s=\mu_d=0$, no viscous
damping and no external loads, for which the total mechanical energy is an invariant of the exact
flow. The endpoint position- and velocity-level constraint norms stay at $1.1\times10^{-13}$ and
$9.5\times10^{-14}$ in both variants, the acceleration-level norms at the stages below
$2\times10^{-14}$, and the quaternion unit error at $2\times10^{-16}$; the relative energy error of
the conservative variant is $4.3\times10^{-12}$ at every step (the exact Gauss step is symplectic
on the reduced problem, and the endpoint closure perturbs it by $O(h^7)$), while the frictional
variant dissipates $37\%$ of its initial energy over the branch. The Newton count is $5$ per step
without friction and $5.7$ with it.

\paperfig{\figpath/Figure_9_e6_regular_branch.png}
{Chain over the regular branch $[0,0.5]$ at $h=0.01$: relative energy error of the conservative
variant (left) and endpoint constraint norms of both variants (right; the velocity-level closure is
applied only when the raw defect exceeds $10^{-13}$, hence the saw-tooth).}
{fig:e6}

\section{Limitations}
\label{sec:limitations}

\begin{enumerate}
  \item \textbf{Scope of the identity.} Lemma~\ref{lem:exact-stage-identity} is stated for the open
  chain, where the joint coordinates form a global chart. For closed loops the same rows are solved,
  but the identity holds only on local charts of independent coordinates, and the theorem does not
  by itself turn the roundoff-level constraint residuals observed on the four-link and slider-crank
  mechanisms into a trajectory-error bound; that needs a stability constant for the mechanism at
  hand.
  \item \textbf{Regular branch.} The frozen sliding direction of the chain's second pair is a
  modelling choice that is regular only while $n_1\cdot a_1(q)\neq0$; on the benchmark initial state
  this holds up to $t\approx0.6$, and all chain experiments stop at $T=0.5$. A pair whose sliding
  direction follows the axis would remove the limit at the cost of a different (still smooth) stage
  system; the analysis applies unchanged.
  \item \textbf{Sharp friction.} The theorem's constants and its threshold $h_0$ are governed by the
  curvature of the regularized friction law. For Stribeck velocities below about $0.1$ the observed
  order is reduced on coarse grids and recovers only for $h\lesssim3\times10^{-3}$
  (Section~\ref{sec:numerical}). Set-valued Coulomb friction is outside the smoothness hypothesis.
  \item \textbf{Solver hypothesis.} Lemma~\ref{lem:newton-envelope} covers predictors within $O(h)$
  of the Gauss stage. The implemented predictor is not of that kind, and the experiments use a fixed
  Newton tolerance rather than the $h^7$-scaled rule, so on the finest grids they satisfy (H3) only
  with a large constant $c_\eta$; that the accepted iterates are nevertheless the Gauss roots to
  within $10^{-11}$ is checked a posteriori, not guaranteed a priori.
  \item \textbf{Cost.} The reference implementation forms and factorizes a dense
  $132\times132$ Jacobian by automatic differentiation in Python; wall-clock times are dominated by
  interpreter overhead and are reported only as indications. A production implementation would
  exploit the stage-block structure and the sparsity of the constraint rows; no scalability study
  in the number of bodies is reported.
  \item \textbf{Comparisons.} The time finite-element method of~\cite{chaturvedi2026tfe} has not
  been implemented; the comparison with it is formal (order six against order five). The public
  absolute-coordinate codes are run as published, with their own tolerances, and are compared on the
  double pendulum only, the one benchmark mechanism whose motion is not prescribed.
\end{enumerate}

\section{Conclusions}
\label{sec:conclusions}

The \method{} step solves a single square stage system that couples three-stage Gauss collocation
of the joint coordinates, the lower-pair constraints at position, velocity and acceleration level,
and the Newton--Euler balance, and reconstructs the rotational endpoint through the exponential
map. The stage system is exactly reduced Gauss collocation lifted to absolute coordinates: the
lifted reduced Gauss stage satisfies all $132$ rows, and every root of the non-dynamic rows on the
regular branch is such a lift. The sixth order of the method is therefore inherited from Gauss
collocation, and the implementation-side perturbations, endpoint closure and inexact Newton, are
$O(h^7)$ under a smoothness hypothesis alone; the interfaces through which the implementation enters
are derived for small steps, with a threshold that is set by the friction law. The algebra and the
perturbation chain are machine-checked.

The experiments confirm the picture at every point where it can be tested: sixth order in position,
orientation, velocity and angular velocity over the regular branch of a frictional chain, with
constraint norms at the $10^{-13}$ level and energy conservation of the frictionless variant to
$10^{-12}$; orders two, four and six for the one-, two- and three-stage members with the same
stage system at the same Newton work; sixth order on the ASME double pendulum and roundoff-level
reproduction of the three driven mechanisms; and a controlled loss of order when the friction law is
sharpened, in the direction the analysis predicts. Against the public absolute-coordinate codes on
the double pendulum over the public horizon, the method reaches errors below $10^{-10}$ at step
sizes where those codes still have errors of order one.

Three directions follow. The identity suggests that other stage-level constraint hierarchies
(Lobatto, Radau) inherit their reduced-problem order in the same way, which would give
$L$-stable members for stiff joints. The $O(h)$ predictor of Lemma~\ref{lem:newton-envelope}
would remove the last a-posteriori element of the analysis at no cost. And a sparse, block-structured
implementation would turn the per-step cost into one comparable to that of the low-order
absolute-coordinate schemes, at which point the work/precision advantage measured here in Newton
iterations would appear in wall-clock time as well.

\section*{CRediT authorship contribution statement}

Jingquan Wang: Conceptualization, Methodology, Software, Validation, Formal analysis, Writing --
original draft, Writing -- review and editing.

\section*{Declaration of competing interest}

The author declares no competing interests relevant to this work.

\section*{Funding}

No external funding was received for this research.

\section*{Data availability}

The implementation, the experiment scripts that generate every table and figure of
Section~\ref{sec:numerical}, the result files, and the Lean~4 development are available in the
repository described in \ref{app:repro}.

\section*{Declaration of generative AI and AI-assisted technologies in the writing process}

During preparation of this work, the author used an AI coding assistant for drafting support,
consistency checking and script editing. The author reviewed and edited the content and takes full
responsibility for the published article.

\appendix

\section{Lower-pair constraint rows of the benchmark mechanisms}
\label{app:lower-pair}

The four benchmark mechanisms use the scalar constraint rows of the public suite
\cite{kissel2022absolute}. Let body $i$ have centre position $r_i$, rotation matrix $A_i\in SO(3)$,
translational velocity $v_i$, angular velocity $\omega_i$, translational acceleration $a_i$ and
angular acceleration $\alpha_i$. A marked point and a marked direction fixed in body $i$ are
$x_i=r_i+A_i s_i$ and $u_i=A_i a_i$. The row families are
\begin{align}
  C_{\mathrm{CD},k}(q,t) &= e_k^T\left(x_j-x_i\right)-d_k(t), &
  C_{\mathrm{DP1}}(q,t) &= u_i^T u_j-c(t),\\
  C_{\mathrm{DP2}}(q,t) &= u_i^T\left(x_j-x_i\right)-d(t), &
  C_{\mathrm{D}}(q,t) &= \left\|x_j-x_i\right\|^2-\ell^2(t),
  \label{eq:lower-pair-row-families}
\end{align}
where CD is a coordinate difference, DP1 constrains two body-fixed directions, DP2 constrains a
body-fixed direction against a point difference, and D is a distance row. Revolute joints are
assembled from CD and DP1 rows; the slider-crank prismatic joint additionally uses DP2 and distance
rows; drivers prescribe $c(t)$ of a DP1 row. With
$\dot x_i = v_i+\omega_i\times A_i s_i$, $\dot u_i = \omega_i\times u_i$,
$\ddot x_i = a_i+\alpha_i\times A_i s_i+\omega_i\times(\omega_i\times A_i s_i)$,
$\ddot u_i = \alpha_i\times u_i+\omega_i\times(\omega_i\times u_i)$, $\Delta x=x_j-x_i$ and
its derivatives, the velocity- and acceleration-level rows enforced at every stage are
\begin{align}
  \dot C_{\mathrm{CD},k} &= e_k^T\Delta\dot x-\dot d_k, &
  \ddot C_{\mathrm{CD},k} &= e_k^T\Delta\ddot x-\ddot d_k,\\
  \dot C_{\mathrm{DP1}} &= \dot u_i^T u_j+u_i^T\dot u_j-\dot c, &
  \ddot C_{\mathrm{DP1}} &= \ddot u_i^T u_j+2\dot u_i^T\dot u_j+u_i^T\ddot u_j-\ddot c,\\
  \dot C_{\mathrm{DP2}} &= \dot u_i^T\Delta x+u_i^T\Delta\dot x-\dot d, &
  \ddot C_{\mathrm{DP2}} &= \ddot u_i^T\Delta x+2\dot u_i^T\Delta\dot x+u_i^T\Delta\ddot x-\ddot d,\\
  \dot C_{\mathrm{D}} &= 2\Delta x^T\Delta\dot x-2\ell\dot\ell, &
  \ddot C_{\mathrm{D}} &= 2\Delta\dot x^T\Delta\dot x+2\Delta x^T\Delta\ddot x-2(\dot\ell^2+\ell\ddot\ell).
  \label{eq:lower-pair-acceleration-rows}
\end{align}

\paragraph{Mechanism data}
All bodies are steel bars of square cross-section $0.05$~m and density $7800$~kg/m$^3$ unless
stated otherwise; $e_x,e_y,e_z$ are the coordinate directions and body-fixed vectors are given in
the body frame.
\emph{Single pendulum}: one bar of length $4$~m (mass $78$~kg,
$J=\operatorname{diag}(0.0325,104.016,104.016)$~kg\,m$^2$), ground revolute joint with CD rows
$s_1=(-2,0,0)$, $s_0=0$ and DP1 rows $(e_x,e_x)$, $(e_y,e_x)$; driven DP1 row $(e_y,-e_z)$ with
$c(t)=\cos\theta(t)$, $\theta(t)=\pi/2+(\pi/4)\cos 2t$.
\emph{Double pendulum}: bars of length $4$ and $2$~m (masses $78$ and $39$~kg), ground revolute
joint as above, inter-link revolute joint with CD rows $s_1=(2,0,0)$, $s_2=(-1,0,0)$ and DP1 rows
on body 2 against ground $(e_x,e_x)$, $(e_y,e_x)$; released from the horizontal configuration
$(\theta_1,\theta_2)(0)=(-\pi/2,0)$ at rest; no friction, no drive.
\emph{Four-link loop}: rotor, link 2 and link 3 with masses $2,1,1$~kg and inertia diagonals
$(4,2,0)$, $(12.4,0.01,0)$, $(4.54,0.01,0)$~kg\,m$^2$; revolute joint A (body 1--ground, CD
$s_1=s_0=0$, DP1 $(e_x,e_x)$, $(e_y,e_x)$), revolute joint D (body 3--ground, CD $s_3=(0,3.7,0)$,
$s_0=(-4,-8.5,0)$, DP1 $((0,0,1),e_y)$, $(e_y,e_y)$), spherical joint (bodies 3--2, CD
$s_3=(0,-3.7,0)$, $s_2=(0,-6.1,0)$), universal joint (bodies 1--2, CD $s_1=(-2,0,0)$,
$s_2=(0,6.1,0)$, DP1 $((0,-0.662,0.75),e_z)$); driving DP1 row $((-1,0,0),e_y)$ with
$c(t)=\cos(\pi t+\pi/2)$.
\emph{Slider-crank}: crank, rod and slider with masses $0.12,0.5,2$~kg and inertia diagonals
$(10^{-4},10^{-5},10^{-4})$, $(4\times10^{-3},4\times10^{-4},4\times10^{-3})$,
$(10^{-4},10^{-4},10^{-4})$~kg\,m$^2$; revolute joint A (ground--crank, CD $s_4=(0,0.1,0.12)$,
$s_1=0$, DP1 $(e_y,e_x)$, $(e_z,e_x)$), spherical joint B (rod--crank, CD $s_2=(-0.15,0,0)$,
$s_1=(0,0.08,0)$), joint C (rod--slider, DP1 $(e_y,e_z)$, two DP2 rows from rod point
$(0.15,0,0)$ to slider point $0$ along $e_y$, $e_z$), translational joint D (slider--ground, DP1
$(-e_x,e_x)$, $(e_y,e_x)$, $(-e_x,e_y)$, DP2 rows from slider point $0$ to ground point
$(0.2,0,0)$ along $-e_x$, $e_y$), distance row between rod point $(0.15,0,0)$ and slider point
$(0,0,1)$ with $\ell=1$; driving DP1 row $(e_y,e_y)$ with $c(t)=\cos(-2\pi t+\pi/2)$.

\section{Lean development}
\label{app:lean}

The Lean~4 development (Mathlib pinned by its manifest) machine-checks the algebraic and
perturbation lemmas of Section~\ref{sec:order}: $57$ theorems, no open goals, and an exhaustive
axiom audit showing that every theorem depends only on the standard axioms
\nolinkurl{propext}, \nolinkurl{Classical.choice} and \nolinkurl{Quot.sound}. It builds with
\texttt{lake build}; the script \nolinkurl{scripts/check.sh} runs the build, the axiom audit and
a lint pass. Table~\ref{tab:lean-development} maps the statements of the paper to the theorems that
check them. The row identities are proved for a transcription of the implemented residual into
Lean; the perturbation lemmas are proved in the abstract normed-space setting with the constants of
the paper. Butcher's order theorem and the Gauss local-error estimate are not formalized and enter
as cited results.

\begin{table}[!htb]
\centering
\caption{Statements of the paper and the Lean theorems that check them.}
\label{tab:lean-development}
\footnotesize
\setlength{\tabcolsep}{4pt}
\begin{tabular}{P{0.40\linewidth}P{0.52\linewidth}}
\toprule
Statement & Lean theorem \\
\midrule
Lemma~\ref{lem:exact-stage-identity}, non-dynamic rows (both directions) &
\nolinkurl{FullVA.Transition.nondynamic_rows_iff} \\
Lemma~\ref{lem:exact-stage-identity}, Newton--Euler rows &
\nolinkurl{NewtonEuler.dynamic_rows_vanish} \\
Lemma~\ref{lem:p2-from-p1}, bound Eq.~\eqref{eq:p2-perturbation-bound} &
\nolinkurl{uniform_inverse_of_perturbation} \\
Lemma~\ref{lem:stage-residual-defect}, root uniqueness &
\nolinkurl{root_unique_of_linearization} \\
Lemma~\ref{lem:gauss-endpoint-defect}, bound Eq.~\eqref{eq:gauss-quadrature-defect} &
\nolinkurl{gauss6_step_defect} \\
Butcher conditions $B(6)$, $C(3)$, $D(3)$; failure of $B(7)$ &
\nolinkurl{Gauss6.butcher_order_six_hypotheses}, \nolinkurl{Gauss6.not_B_seven} \\
Lemma~\ref{lem:endpoint-closure} &
\nolinkurl{endpoint_closure_exists}, \nolinkurl{endpoint_correction_bound_h7} \\
Lemma~\ref{lem:inexact-newton} &
\nolinkurl{inexact_newton_output_bound} \\
Lemma~\ref{lem:newton-envelope}, decay Eq.~\eqref{eq:newton-envelope-decay} &
\nolinkurl{simplified_newton_residual_decay} \\
Lemma~\ref{lem:local-global-transfer} &
\nolinkurl{local_to_global}, \nolinkurl{reported_grid_bound} \\
Theorem~\ref{thm:g6fullva-order}, three-term chain and grid bounds &
\nolinkurl{local_defect_bound}, \nolinkurl{conditional_sixth_order_grid_bound} \\
\bottomrule
\end{tabular}
\end{table}

\section{Reproducibility}
\label{app:repro}

The repository accompanying the paper has five top-level folders: \nolinkurl{paper/} (this
manuscript), \nolinkurl{proof/} (the Lean development), \nolinkurl{numerics/} (the implementation
and all experiment scripts), \nolinkurl{external/} (the public benchmark codes as git submodules or
mirrors) and \nolinkurl{validation/} (the audit trail of the development, including the check that
the numbers in this paper equal the values in the result files). The implementation of the step is
the single residual function in \nolinkurl{numerics/v047_cylindrical_chain_pipeline/run_v047.py};
the experiments of Section~\ref{sec:numerical} are the scripts \nolinkurl{e1_convergence.py} to
\nolinkurl{e7_solver_envelope.py} in \nolinkurl{numerics/v049_paper_experiments/}, each of which
writes a CSV and a JSON file under \nolinkurl{results/}; the figures are drawn from those files by
\nolinkurl{make_figures.py}, and \nolinkurl{validation/validate_manuscript.py} checks that every
number in the tables of that section equals the value in its result file. The Lean development is checked by
\nolinkurl{proof/scripts/check.sh}. All runs reported were made on one CPU core; the double
pendulum reference at $h=10^{-4}$ over $T=3$ took $1.6$~h, each of the other scripts between a
minute and about an hour.

\begin{thebibliography}{00}

\bibitem{chaturvedi2026tfe}
E. Chaturvedi, C. Sandu, A. Sandu,
Higher-order integration of index-3 DAE with friction using time finite elements on Lie groups,
Multibody System Dynamics (2026).
doi:10.1007/s11044-026-10153-w.

\bibitem{kim2017higher}
W. Kim, J.N. Reddy,
A new family of higher-order time integration algorithms for the analysis of structural dynamics,
Journal of Applied Mechanics 84 (7) (2017) 071008.
doi:10.1115/1.4036821.

\bibitem{hairer1993solving}
E. Hairer, G. Wanner,
Solving Ordinary Differential Equations II: Stiff and Differential-Algebraic Problems,
second ed., Springer, Berlin, 1996.

\bibitem{hairer1993nonstiff}
E. Hairer, S.P. N{\o}rsett, G. Wanner,
Solving Ordinary Differential Equations I: Nonstiff Problems,
second ed., Springer, Berlin, 1993.

\bibitem{hairer2006geometric}
E. Hairer, C. Lubich, G. Wanner,
Geometric Numerical Integration: Structure-Preserving Algorithms for Ordinary Differential
Equations, second ed., Springer, Berlin, 2006.

\bibitem{butcher2016numerical}
J.C. Butcher,
Numerical Methods for Ordinary Differential Equations, third ed., Wiley, Chichester, 2016.

\bibitem{ascher1998computer}
U.M. Ascher, L.R. Petzold,
Computer Methods for Ordinary Differential Equations and Differential-Algebraic Equations,
SIAM, Philadelphia, 1998.

\bibitem{petzold1986numerical}
L.R. Petzold, P. L\"otstedt,
Numerical solution of nonlinear differential equations with algebraic constraints II: practical
implications,
SIAM Journal on Scientific and Statistical Computing 7 (3) (1986) 720--733.

\bibitem{gear1985automatic}
C.W. Gear, B. Leimkuhler, G.K. Gupta,
Automatic integration of Euler--Lagrange equations with constraints,
Journal of Computational and Applied Mathematics 12--13 (1985) 77--90.
doi:10.1016/0377-0427(85)90008-1.

\bibitem{bauchau2008review}
O.A. Bauchau, A. Laulusa,
Review of contemporary approaches for constraint enforcement in multibody systems,
Journal of Computational and Nonlinear Dynamics 3 (1) (2008) 011005.

\bibitem{crouch1993numerical}
P.E. Crouch, R. Grossman,
Numerical integration of ordinary differential equations on manifolds,
Journal of Nonlinear Science 3 (1993) 1--33.

\bibitem{munthe1998lie}
H. Munthe-Kaas,
Runge--Kutta methods on Lie groups,
BIT Numerical Mathematics 38 (1998) 92--111.

\bibitem{munthe1999high}
H. Munthe-Kaas,
High order Runge--Kutta methods on manifolds,
Applied Numerical Mathematics 29 (1999) 115--127.

\bibitem{iserles2000lie}
A. Iserles, H.Z. Munthe-Kaas, S.P. N{\o}rsett, A. Zanna,
Lie-group methods,
Acta Numerica 9 (2000) 215--365.

\bibitem{celledoni2014lie}
E. Celledoni, H. Marthinsen, B. Owren,
An introduction to Lie group integrators: basics, new developments and applications,
Journal of Computational Physics 257 (2014) 1040--1061.

\bibitem{bourabee2009hamilton}
N. Bou-Rabee, J.E. Marsden,
Hamilton--Pontryagin integrators on Lie groups. Part I: Introduction and structure-preserving
properties,
Foundations of Computational Mathematics 9 (2009) 197--219.

\bibitem{bottasso1998integrating}
C.L. Bottasso, M. Borri,
Integrating finite rotations,
Computer Methods in Applied Mechanics and Engineering 164 (1998) 307--331.

\bibitem{arnold2007convergence}
M. Arnold, O. Br\"uls,
Convergence of the generalized-$\alpha$ scheme for constrained mechanical systems,
Multibody System Dynamics 18 (2007) 185--202.

\bibitem{bruls2010lie}
O. Br\"uls, A. Cardona,
On the use of Lie group time integrators in multibody dynamics,
Journal of Computational and Nonlinear Dynamics 5 (3) (2010) 031002.

\bibitem{bruls2012lie}
O. Br\"uls, A. Cardona, M. Arnold,
Lie group generalized-$\alpha$ time integration of constrained flexible multibody systems,
Mechanism and Machine Theory 48 (2012) 121--137.
doi:10.1016/j.mechmachtheory.2011.07.017.

\bibitem{arnold2015error}
M. Arnold, O. Br\"uls, A. Cardona,
Error analysis of generalized-$\alpha$ Lie group time integration methods for constrained
mechanical systems,
Numerische Mathematik 129 (2015) 149--179.

\bibitem{sonneville2014geometrically}
V. Sonneville, A. Cardona, O. Br\"uls,
Geometrically exact beam finite element formulated on the special Euclidean group $SE(3)$,
Computer Methods in Applied Mechanics and Engineering 268 (2014) 451--474.

\bibitem{wieloch2021bdf}
V. Wieloch, M. Arnold,
BDF integrators for constrained mechanical systems on Lie groups,
Journal of Computational and Applied Mathematics 387 (2021) 112517.
doi:10.1016/j.cam.2019.112517.

\bibitem{terze2016singularity}
Z. Terze, A. M\"uller, D. Zlatar,
Singularity-free time integration of rotational quaternions using non-redundant ordinary
differential equations,
Multibody System Dynamics 38 (2016) 201--225.
doi:10.1007/s11044-016-9518-7.

\bibitem{terze2020aircraft}
Z. Terze, D. Zlatar, V. Pandza,
Aircraft attitude reconstruction via novel quaternion-integration procedure,
Aerospace Science and Technology 97 (2020) 105617.

\bibitem{lunk2006solving}
C. Lunk, B. Simeon,
Solving constrained mechanical systems by the family of Newmark and $\alpha$-methods,
ZAMM -- Journal of Applied Mathematics and Mechanics 86 (10) (2006) 772--784.

\bibitem{jay1996symplectic}
L.O. Jay,
Symplectic partitioned Runge--Kutta methods for constrained Hamiltonian systems,
SIAM Journal on Numerical Analysis 33 (1) (1996) 368--387.
doi:10.1137/0733019.

\bibitem{jay2006specialized}
L.O. Jay,
Specialized Runge--Kutta methods for index 2 differential-algebraic equations,
Mathematics of Computation 75 (254) (2006) 641--654.

\bibitem{jay2007specialized}
L.O. Jay,
Specialized partitioned additive Runge--Kutta methods for systems of overdetermined DAEs with
holonomic constraints,
SIAM Journal on Numerical Analysis 45 (5) (2007) 1814--1842.

\bibitem{taves2020exponential}
J. Taves, A. Kissel, D. Negrut,
On an exponential map approach for rigid body kinematics and dynamics analysis,
Technical Report TR-2020-08, Simulation-Based Engineering Laboratory, University of
Wisconsin--Madison (2020).

\bibitem{kissel2021dwelling}
A. Kissel, D. Negrut, J. Taves,
Dwelling on the connection between SO(3) and rotation matrices in rigid multibody dynamics.
Part 1: Description of an index-3 DAE solution approach,
in: Proceedings of the ASME 2021 International Design Engineering Technical Conferences and
Computers and Information in Engineering Conference, IDETC-CIE2021, 2021.

\bibitem{kissel2022absolute}
A. Kissel, D. Negrut, J. Taves,
Constrained multibody kinematics and dynamics in absolute coordinates: a discussion of three
approaches to representing rigid body rotation,
Journal of Computational and Nonlinear Dynamics 17 (10) (2022) 101008.

\bibitem{fang2023halfimplicit}
L. Fang, A. Kissel, R. Zhang, D. Negrut,
On the use of half-implicit numerical integration in multibody dynamics,
Journal of Computational and Nonlinear Dynamics 18 (1) (2023) 014501.
doi:10.1115/1.4056183.

\bibitem{kissel2024velocitypartitioning}
A. Kissel, L. Bakke, D. Negrut,
Reducing the constrained multibody dynamics problem to the solution of a system of ordinary
differential equations via velocity partitioning and Lie group integration,
Journal of Computational and Nonlinear Dynamics 19 (7) (2024).
doi:10.1115/1.4065254.

\bibitem{nikravesh1988computer}
P.E. Nikravesh,
Computer-Aided Analysis of Mechanical Systems, Prentice-Hall, Englewood Cliffs, 1988.

\bibitem{shabana2020dynamics}
A.A. Shabana,
Dynamics of Multibody Systems, fifth ed., Cambridge University Press, Cambridge, 2020.

\bibitem{jain2011robot}
A. Jain,
Robot and Multibody Dynamics: Analysis and Algorithms, Springer, New York, 2011.

\bibitem{marsden2001discrete}
J.E. Marsden, M. West,
Discrete mechanics and variational integrators,
Acta Numerica 10 (2001) 357--514.
doi:10.1017/S096249290100006X.

\bibitem{leyendecker2008variational}
S. Leyendecker, J.E. Marsden, M. Ortiz,
Variational integrators for constrained dynamical systems,
ZAMM -- Journal of Applied Mathematics and Mechanics 88 (9) (2008) 677--708.
doi:10.1002/zamm.200700173.

\bibitem{leok2012prolongation}
M. Leok, T. Shingel,
Prolongation--collocation variational integrators,
IMA Journal of Numerical Analysis 32 (3) (2012) 1194--1216.
doi:10.1093/imanum/drr042.

\bibitem{betsch2002conservation}
P. Betsch, P. Steinmann,
Conservation properties of a time FE method. Part III: Mechanical systems with holonomic
constraints,
International Journal for Numerical Methods in Engineering 53 (10) (2002) 2271--2304.

\bibitem{browmcphee2016friction}
P. Brown, J. McPhee,
A continuous velocity-based friction model for dynamics and control with physically meaningful
parameters,
Journal of Computational and Nonlinear Dynamics 11 (5) (2016) 054502.

\bibitem{pennestri2016review}
E. Pennestr\`{\i}, V. Rossi, P. Salvini, P.P. Valentini,
Review and comparison of dry friction force models,
Nonlinear Dynamics 83 (2016) 1785--1801.

\bibitem{stewart1996implicit}
D.E. Stewart, J.C. Trinkle,
An implicit time-stepping scheme for rigid body dynamics with inelastic collisions and Coulomb
friction,
International Journal for Numerical Methods in Engineering 39 (15) (1996) 2673--2691.

\bibitem{anitescu1997formulating}
M. Anitescu, F.A. Potra,
Formulating dynamic multi-rigid-body contact problems with friction as solvable linear
complementarity problems,
Nonlinear Dynamics 14 (3) (1997) 231--247.
doi:10.1023/A:1008292328909.

\bibitem{tasora2011matrixfree}
A. Tasora, M. Anitescu,
A matrix-free cone complementarity approach for solving large-scale, nonsmooth, rigid body
dynamics,
Computer Methods in Applied Mechanics and Engineering 200 (2011) 439--453.
doi:10.1016/j.cma.2010.06.030.

\bibitem{borri1988rigid}
M. Borri, S.N. Atluri,
Time-finite element method for the constrained dynamics of a rigid body,
in: Computational Mechanics, Springer, Berlin, 1988.

\bibitem{hughes1988spacetime}
T.J.R. Hughes, G.M. Hulbert,
Space-time finite element formulations for elastodynamics: formulations and error estimates,
Computer Methods in Applied Mechanics and Engineering 66 (1988) 339--363.

\bibitem{hulbert1992time}
G.M. Hulbert,
Time finite element methods for structural dynamics,
International Journal for Numerical Methods in Engineering 33 (1992) 307--331.

\bibitem{jog2016robust}
C.S. Jog, M. Agrawal, A. Nandy,
The time finite element as a robust general scheme for solving nonlinear dynamic equations
including chaotic systems,
Applied Mathematics and Computation 279 (2016) 43--61.

\bibitem{bauchau2024time}
O.A. Bauchau,
The finite element method in time for multibody dynamics,
Journal of Computational and Nonlinear Dynamics 19 (7) (2024) 071001.
doi:10.1115/1.4063953.

\end{thebibliography}

\end{document}
